\section{Progression-free sets of integers and of ideals}\label{sec:gpi}
% Displayed path of a delivered file (a long path cannot be broken inside a line)
\providecommand{\gpifile}[2][]{\begin{center}\footnotesize\artifact{#2}#1\end{center}}


This section applies the exponent principle of the source construction, described at the
start of \cref{sec:gpp}, to sets of integers and of ideals without three-term geometric
progressions: $ac=b^2$ holds exactly when $v_p(a)+v_p(c)=2v_p(b)$ at every prime $p$. It
answers density questions of McNew \cite{McNew2015Sets} over the integers and of Best, Huan,
McNew, Miller, Powell, Tor and Weinstein \cite{Best2017Geometric} over number fields, and the
corresponding question of Brown and Gordon \cite{BrownGordon1996Sequences} for geometric
progressions of any length. Every result of this section is unconditional; none uses
\cref{src:thm:dgrh} or any other statement about zeros. Finite computations enter only
through certificates in exact integer or rational arithmetic, which are cited with the result
they support.

\subsection{Conventions and the questions}\label{gpi:sec:setting}
{\raggedright\sources{grh_transfers_density_gaps_and_zero_tests; mcnew_density_and_grh_transfers; grh_ideal_densities_and_transfer_conditions_2026-11-26; grh_prime_patterns_and_branch_precision; prime_observations_and_progression_free_sets}}

Let $K$ be a number field of degree $n=[K:\Q]$ and let $\mathcal I_K$ be the set of its
nonzero integral ideals. Put
\[
 A_K(X)=\#\set{\mathfrak a\in\mathcal I_K:\Nm\mathfrak a\le X},\qquad
 w_{\mathfrak p}=(\Nm\mathfrak p)^{-1}\quad(\mathfrak p\text{ a prime ideal}).
\]
The classical ideal-counting theorem gives $A_K(X)=\kappa_KX+o(X)$ with $\kappa_K>0$
\cite{Murty2007Counting}. Consequently $A_K(t)\ll_Kt$ for all $t\ge1$, and for a fixed ideal
$\mathfrak d$ the multiples of $\mathfrak d$ of norm at most $X$ number
$A_K(X/\Nm\mathfrak d)$, which is $(\Nm\mathfrak d)^{-1}A_K(X)(1+o(1))$ and is
$O_K(X/\Nm\mathfrak d)$ uniformly in $X$ and $\mathfrak d$. The \emph{natural density} of
$E\subseteq\mathcal I_K$ is
\[
 d(E)=\lim_{X\to\infty}\frac{\#\set{\mathfrak a\in E:\Nm\mathfrak a\le X}}{A_K(X)}
\]
when the limit exists, and its \emph{upper density} $\overline d(E)$ is the corresponding
limsup. The count is of ideals; for $K=\Q$ the ideals are the positive integers and these are
the usual densities of subsets of $\N=\set{1,2,\dots}$.

We use two conventions for the forbidden triples. In the \emph{integral} convention they are
$\mathfrak a,\mathfrak a\mathfrak r,\mathfrak a\mathfrak r^2$ with an integral ideal
$\mathfrak r\ne\OO_K$. In the \emph{fractional} convention $\mathfrak r\ne\OO_K$ may be any
nonzero fractional ideal such that all three terms are integral; this includes nontrivial
ratios of norm one. In valuation coordinates
$\mathfrak a\mapsto(v_{\mathfrak p}(\mathfrak a))_{\mathfrak p}$ the forbidden triples are the
nonconstant three-term arithmetic progressions of exponent vectors with a nonnegative step,
respectively with an arbitrary integral step. For $K=\Q$ the integral convention forbids
$a,ar,ar^2$ with an integer $r>1$, and the fractional convention forbids them for every
rational $r>1$, that is, it forbids every triple of distinct positive integers $a,b,c$ with
$ac=b^2$. Thus $4,6,9$ is forbidden only by the second rule. A set is \emph{admissible} (for
the convention under discussion) if it contains no forbidden triple. A set that is admissible
for the fractional convention is admissible for the integral one.

For a fixed convention let $b_K$ be the supremum of the natural densities of the admissible
sets that have a natural density, and $u_K$ the supremum of the upper densities of admissible
sets. For $K=\Q$ we use McNew's notation:
\[
 \alpha,\ \overline\alpha\quad\text{(rational ratios)},\qquad
 \beta,\ \overline\beta\quad\text{(integer ratios)}
\]
for the natural-density suprema $\alpha,\beta$ and the upper-density suprema
$\overline\alpha,\overline\beta$. Trivially $\alpha\le\overline\alpha$,
$\beta\le\overline\beta$, $\alpha\le\beta$ and $\overline\alpha\le\overline\beta$. Results
about one of these four numbers do not transfer to another: an upper bound for natural
densities is not an upper bound for upper densities, and a set that avoids integer ratios
need not avoid rational ones.

\paragraph{The greedy set.}
Let $\mathcal A$ be the set of nonnegative integers whose ternary digits are all $0$ or $1$;
it contains no nonconstant three-term arithmetic progression, by the carry-free argument
of \cref{gpp:sec:conventions}. Note also that
$\mathcal A=\set{0,1}+3\mathcal A=\set{0,1,3,4}+9\mathcal A$ with unique representations. The
\emph{greedy set} is
\[
 \mathcal G_K=\set{\mathfrak a\in\mathcal I_K:v_{\mathfrak p}(\mathfrak a)\in\mathcal A
 \text{ for every prime ideal }\mathfrak p}.
\]
It is admissible for the fractional convention. For $K=\Q$ it is Rankin's set
$R=\mathcal G_\Q$ \cite{Rankin1961Sets}, and by Lemma~\ref{gpi:lem:product} below
\[
 d(R)=\prod_p\Bigl[(1-p^{-1})\prod_{j\ge0}(1+p^{-3^j})\Bigr]\approx0.71974 .
\]

\paragraph{The questions.}
McNew \cite[Section~7]{McNew2015Sets} asked whether $\alpha<\overline\alpha$ (Question~7.1),
whether $\beta<\overline\beta$ (Question~7.3) and whether $\alpha<\beta$ (Question~7.4), and
conjectured that $\alpha=d(R)$ (Conjecture~7.2). Best et al.\ ask how far nongreedy sets of
ideals in a quadratic field can improve on the greedy set, in particular for natural density
with integral ratios (their Question~5.1) and with fractional ratios (their Question~5.4)
\cite[Section~5]{Best2017Geometric}. The answers obtained in this section are as follows.
\begin{itemize}
\item Questions~7.1 and~7.3: $\overline\alpha-\alpha>0.002836$ and
 $\overline\beta-\beta>0.069218$ (\cref{gpi:thm:integer-gap}). More generally $u_K>b_K$ in
 every number field and for both conventions (\cref{gpi:thm:ideal-separation}).
\item Conjecture~7.2 is false: $\alpha>d(R)$. A change of the lowest ternary digit of the
 exponents at finitely many primes suffices (\cref{gpi:thm:onedigit}), and an explicit set
 supported on the primes up to $1427$ has $d(E)-d(R)>10^{-596}$ (\cref{gpi:thm:explicit}).
 On the other hand Rankin's rule is optimal on every block of at most four primes, so that
 $\alpha<0.740912$ (Theorem~\ref{gpi:thm:fourprimes} and Corollary~\ref{gpi:cor:restricted}), and every density rule that
 couples only boundedly many primes at a time loses a positive amount against $\alpha$
 (\cref{gpi:thm:blocks}). The suprema $\alpha$ and $\beta$ are computable and are attained
 (\cref{gpi:thm:attain}).
\item For progressions of every length $k\ge3$, with rational ratios, the greedy set is surpassed
 in natural density, the gain can be prescribed independently in the reduced residue classes of any
 modulus, and every bound on the number of jointly tested primes costs a positive amount of density
 (\cref{gpi:thm:klength,gpi:thm:kresidues,gpi:thm:blocksk}); this answers the question of Brown and
 Gordon \cite{BrownGordon1996Sequences} recalled in \cref{gpi:sec:klength}.
\item Question~7.4: $\beta-\alpha>0.0001351$ (\cref{gpi:thm:beta-alpha}).
\item In every number field the greedy ideals are surpassed in natural density, even when all
 fractional-ideal ratios are forbidden (\cref{gpi:thm:onedigit}); for integral ratios in
 $\Q(\sqrt{-23})$ the gain exceeds $0.0011957$ (Proposition~\ref{gpi:prop:Q23gain}).
\item For the upper density with rational ratios, $0.7365954<\overline\alpha<0.7716011$
 (\cref{gpi:thm:upper-alpha}).
\end{itemize}
The same questions over $\F_q[x]$, with the same arguments, are in \cref{gpi:sec:poly}, a
subsection of \cref{sec:gpp}. The following table lists the results by question and ring.
\begin{center}\small
\begin{tabular}{@{}>{\raggedright}p{2.9cm}>{\raggedright}p{3.6cm}>{\raggedright}p{2.9cm}>{\raggedright\arraybackslash}p{4.4cm}@{}}
\toprule
Question & Integers & Number fields & $\F_q[x]$ and finite rings\\
\midrule
Natural against upper density & \cref{gpi:thm:integer-gap} &
 \cref{gpi:thm:ideal-separation} &
 \cref{gpi:thm:poly-gap,gpi:rem:poly-ratio}; finite rings: \cref{gpp:thm:gaprings}\\
Computability and attainment & \cref{gpi:thm:attain}; every length:
 \cref{gpi:thm:blocksk} & computability: \cref{gpi:thm:ideal-separation}; attainment: -- &
 \cref{gpp:thm:supremum,gpi:thm:poly-natural}\\
The greedy set is not optimal &
 \cref{gpi:thm:onedigit,gpi:thm:explicit,gpi:thm:klength} &
 \cref{gpi:thm:onedigit,gpi:prop:cube} &
 \cref{gpp:thm:natural,gpp:thm:rational,gpi:prop:poly-greedy}\\
Where the greedy rule is optimal &
 \cref{gpi:prop:light,gpi:thm:fourprimes,gpi:cor:restricted,gpi:thm:blocks} & -- & --\\
The degree-deficit benchmark $m_q$ & -- & -- &
 \cref{gpp:thm:counterexample,gpp:thm:twoprime,gpp:thm:largeq}\\
\bottomrule
\end{tabular}
\end{center}

\subsection{Product sets and the two smooth-factor optimizations}\label{gpi:sec:local}
{\raggedright\sources{grh_transfers_density_gaps_and_zero_tests; grh_ideal_densities_and_transfer_conditions_2026-11-26; mcnew_counterexample_and_grh_transfers; grh_prime_patterns_and_branch_precision; grh_transfer_advances_2026-11-26}}

Let $S$ be a finite set of prime ideals. An ideal is \emph{$S$-smooth} if all its prime
factors lie in $S$ and \emph{$S$-rough} if none does; every ideal factors uniquely as
$\mathfrak a=\mathfrak g\mathfrak h$ with $\mathfrak g$ rough and $\mathfrak h$ smooth. For a
set $F$ of $S$-smooth ideals put
\[
 W(F)=\sum_{\mathfrak h\in F}\frac1{\Nm\mathfrak h},\qquad
 c_S=\prod_{\mathfrak p\in S}(1-w_{\mathfrak p}),
\]
so that the set of all $S$-smooth ideals has $W=1/c_S$, by multiplying geometric series. For
$K=\Q$ and a real $Y$ we write $S_Y$ for the set of positive integers all of whose prime
factors are at most $Y$.

\begin{lemma}[Density of a product set]\label{gpi:lem:product}
Let $S$ be a finite set of prime ideals, let $F$ be a set of $S$-smooth ideals, and for every
prime $\mathfrak p\notin S$ let $A_{\mathfrak p}$ be a set of nonnegative integers containing
$0$ and $1$. Let $E$ consist of the ideals $\mathfrak g\mathfrak h$ with $\mathfrak h\in F$
and $\mathfrak g$ an $S$-rough ideal such that $v_{\mathfrak p}(\mathfrak g)\in A_{\mathfrak p}$
for every $\mathfrak p\notin S$. Then $E$ has the natural density
\begin{equation}\label{gpi:eq:product}
 d(E)=c_S\,W(F)\prod_{\mathfrak p\notin S}\lambda_{\mathfrak p},\qquad
 \lambda_{\mathfrak p}=(1-w_{\mathfrak p})\sum_{e\in A_{\mathfrak p}}w_{\mathfrak p}^{e}
 \ \ge\ 1-w_{\mathfrak p}^2 ,
\end{equation}
and the infinite product converges to a positive number. If every $A_{\mathfrak p}$ contains
no nonconstant three-term arithmetic progression and $F$ is admissible, then $E$ is
admissible for the same convention.
\end{lemma}

\begin{proof}
Counting multiples of products of prime powers and applying finite inclusion--exclusion
shows that, for a finite set $T$ of primes and prescribed exponents $(e_{\mathfrak p})$, the
ideals with $v_{\mathfrak p}=e_{\mathfrak p}$ for all $\mathfrak p\in T$ have density
$\prod_{\mathfrak p\in T}(1-w_{\mathfrak p})w_{\mathfrak p}^{e_{\mathfrak p}}$. The ideals with
$v_{\mathfrak p}\ge L$ for some $\mathfrak p\in T$ have upper density
$O_K(\sum_{\mathfrak p\in T}w_{\mathfrak p}^L)$, which tends to zero as $L\to\infty$. Hence
any condition on the valuations at the primes of $T$ defines a set whose density is the
corresponding sum of products of the independent geometric laws.

Let $H$ be the set of rough ideals occurring in the definition of $E$. Imposing its
conditions only at the primes of norm at most $Y$ changes the set only inside the ideals
divisible by $\mathfrak p^2$ for some $\mathfrak p$ with $\Nm\mathfrak p>Y$, because
$0,1\in A_{\mathfrak p}$. The number of such ideals of norm at most $X$ is at most
$\sum_{\Nm\mathfrak p>Y}A_K(X/\Nm\mathfrak p^2)$, so their upper density is
$O_K(\sum_{\Nm\mathfrak p>Y}w_{\mathfrak p}^2)$, which tends to zero as $Y\to\infty$ since at
most $n$ primes have a given norm. Therefore $d(H)=c_H:=c_S\prod_{\mathfrak p\notin S}
\lambda_{\mathfrak p}$. The product is positive because
$\lambda_{\mathfrak p}\ge(1-w_{\mathfrak p})(1+w_{\mathfrak p})$ and
$\sum_{\mathfrak p}w_{\mathfrak p}^2<\infty$. The same argument gives
$d(\mathfrak hH)=c_H/\Nm\mathfrak h$ for every smooth $\mathfrak h$. By unique factorization
the sets $\mathfrak hH$ are disjoint for distinct $\mathfrak h$, and the union of those with
$\Nm\mathfrak h>T$ has upper density $O_K(\sum_{\Nm\mathfrak h>T}1/\Nm\mathfrak h)$, which
tends to zero as $T\to\infty$. Hence $d(E)=c_HW(F)$.

In a forbidden triple inside $E$ the valuations at each $\mathfrak p\notin S$ form an
arithmetic progression in $A_{\mathfrak p}$, so they are constant. The three rough parts
agree, and cancelling the common rough part leaves a forbidden triple in $F$ with the same
ratio.
\end{proof}

Two choices recur: $A_{\mathfrak p}=\mathcal A$ (the greedy rule outside $S$), for which
$\lambda_{\mathfrak p}=(1-w_{\mathfrak p})\prod_{j\ge0}(1+w_{\mathfrak p}^{3^j})$, and
$A_{\mathfrak p}=\set{0,1}$ (squarefree rough parts), for which
$\lambda_{\mathfrak p}=1-w_{\mathfrak p}^2$. Taking $S=\emptyset$ in the first gives
$d(\mathcal G_K)=\prod_{\mathfrak p}(1-w_{\mathfrak p})\sum_{e\in\mathcal A}w_{\mathfrak p}^e$.

\paragraph{The two optimizations.}
Fix a convention and an integer $Y\ge2$, and let $S$ be the set of prime ideals of norm at
most $Y$. Put
\[
 c=c_S,\qquad v=\prod_{\mathfrak q\notin S}\bigl(1-w_{\mathfrak q}^{2}\bigr),
\]
and, for admissible sets $D$ of $S$-smooth ideals and $t\ge1$,
\begin{align*}
 B&=c\,\sup_{D}W(D),\\
 J(t)&=\max\set{\#D:D\text{ admissible},\ \Nm\mathfrak h\le t\text{ for all }\mathfrak h\in D},\\
 U&=c\int_1^\infty J(t)\,\frac{\dd t}{t^2}.
\end{align*}
The number $B$ optimizes one fixed admissible smooth set against the weights
$1/\Nm\mathfrak h$; the number $U$ optimizes separately at every cutoff. If $H_S(t)$ counts the
smooth ideals of norm at most $t$, then $c\int_1^\infty H_S(t)t^{-2}\,\dd t$ equals $c$ times
the reciprocal-norm sum of all smooth ideals, which is $1$; and for a fixed admissible $D$
with counting function $D(t)$
\begin{equation}\label{gpi:eq:tonelli}
 \int_1^\infty D(t)\,\frac{\dd t}{t^2}=W(D),\qquad D(t)\le J(t)\le H_S(t).
\end{equation}
Hence $B\le U\le1$. The supremum defining $B$ is attained: identify subsets of the countable
smooth monoid with their indicator vectors in a product of two-point spaces; admissibility is
a closed condition and the summable weighted objective is continuous.

\begin{lemma}[Smooth-factor bounds]\label{gpi:lem:BU}
For either convention and every $Y$,
\[
 vB\le b_K\le B,\qquad vU\le u_K\le U,\qquad
 1-v\le\sum_{\Nm\mathfrak q>Y}w_{\mathfrak q}^2\le\frac nY .
\]
For $K=\Q$, moreover, every admissible set $E\subseteq\N$ satisfies
\[
 \limsup_{X\to\infty}\frac1{\log X}\sum_{m\in E,\ m\le X}\frac1m\ \le\ B .
\]
\end{lemma}

\begin{proof}
By finite inclusion--exclusion the rough ideals have density $c$, and by
Lemma~\ref{gpi:lem:product} the squarefree rough ideals have density $cv$.

\emph{Upper bound for $b_K$.} Let $E$ be admissible with natural density $\delta$. Partial
summation gives
$\sum_{\mathfrak a\in E,\ \Nm\mathfrak a\le X}1/\Nm\mathfrak a=(\delta\kappa_K+o(1))\log X$.
For each rough $\mathfrak g$ the fiber $D_{\mathfrak g}=\set{\mathfrak h\text{ smooth}:
\mathfrak g\mathfrak h\in E}$ is admissible, since multiplication by $\mathfrak g$ preserves
forbidden triples; no density is assumed to exist within a fiber. Hence
$W(D_{\mathfrak g})\le B/c$ and
\[
 \sum_{\substack{\mathfrak a\in E\\ \Nm\mathfrak a\le X}}\frac1{\Nm\mathfrak a}
 \le\frac Bc\sum_{\substack{\mathfrak g\text{ rough}\\ \Nm\mathfrak g\le X}}\frac1{\Nm\mathfrak g}
 =(B\kappa_K+o(1))\log X ,
\]
again by partial summation. Thus $\delta\le B$. For $K=\Q$ the same inequality, with a limsup
on the left, gives the last assertion of the lemma.

\emph{Upper bound for $u_K$.} For admissible $E$, counting in each rough fiber gives
\[
 \#\set{\mathfrak a\in E:\Nm\mathfrak a\le X}
 \le\sum_{\substack{\mathfrak g\text{ rough}\\ \Nm\mathfrak g\le X}}J(X/\Nm\mathfrak g).
\]
Fix $T>1$. On $[1,T]$ the function $J$ is a step function with finitely many jumps, and on
an interval $(a,b)$ where it is constant the rough ideals with $X/b<\Nm\mathfrak g<X/a$ have
count $(c(1/a-1/b)+o(1))A_K(X)$. So the terms with $X/\Nm\mathfrak g\le T$ contribute
$(c\int_1^TJ(t)t^{-2}\,\dd t+o(1))A_K(X)$. The remaining terms are at most the number of
ideals $\mathfrak g\mathfrak h$ of norm at most $X$ with $\Nm\mathfrak g<X/T$. Those with
$\Nm\mathfrak h>\sqrt T$ number at most
$\sum_{\Nm\mathfrak h>\sqrt T}A_K(X/\Nm\mathfrak h)\ll_KX\sum_{\Nm\mathfrak h>\sqrt T}1/\Nm\mathfrak h$,
the sums being over smooth $\mathfrak h$, and the others have norm less than $X/\sqrt T$.
Both bounds are $o(X)$ as $T\to\infty$, uniformly in $X$. Letting $X\to\infty$ and then
$T\to\infty$ gives $\overline d(E)\le U$.

\emph{Lower bounds.} Retain only squarefree rough parts. By Lemma~\ref{gpi:lem:product} with
$A_{\mathfrak q}=\set{0,1}$, using one admissible smooth set $D$ in every fiber gives an
admissible set of natural density $cvW(D)$; the supremum over $D$ is $vB$. For the upper
density choose integers $X_j$ with $X_0=0$ and $X_j/X_{j-1}\to\infty$. For every squarefree
rough $\mathfrak g$ with $X_{j-1}<\Nm\mathfrak g\le X_j$ permanently assign a maximizing
admissible smooth set of ideals of norm at most $X_j/\Nm\mathfrak g$, and let $E$ be the union
over all $\mathfrak g$ and $j$ of the assigned sets multiplied by $\mathfrak g$. In a forbidden
triple inside $E$ the valuations outside $S$ form arithmetic progressions in $\set{0,1}$, so
the three terms lie in one fiber, where the assigned set is admissible. Thus $E$ is
admissible. At the cutoff $X=X_j$ and for fixed $T$, the fibers with
$X/T<\Nm\mathfrak g\le X$ were all assigned at stage $j$ once $j$ is large, and by the
step-function count above they contribute $(cv\int_1^TJ(t)t^{-2}\,\dd t+o(1))A_K(X)$. Letting
$T\to\infty$ gives $\overline d(E)\ge vU$.

\emph{The tail.} A prime ideal of norm $m$ lies over the unique rational prime dividing $m$,
and at most $n$ prime ideals lie over it. Hence
$1-v\le\sum_{\Nm\mathfrak q>Y}w_{\mathfrak q}^2\le n\sum_{m>Y}m^{-2}\le n/Y$.
\end{proof}

For $K=\Q$ the sharper form $1-v\le\sum_{p>Y}p^{-2}$ of the tail bound will be used. The lemma
has the following immediate consequence, on which all the separation results rest:
\begin{equation}\label{gpi:eq:gapframe}
 u_K-b_K\ \ge\ vU-B\ \ge\ (U-B)-(1-v),\qquad
 U-B=c\,\inf_D\int_1^\infty\bigl(J(t)-D(t)\bigr)\frac{\dd t}{t^2},
\end{equation}
where $D$ ranges over the admissible smooth sets and the second inequality uses $U\le1$. A
lower bound for $U-B$ is thus a statement that no single admissible set is optimal at every
cutoff.

\subsection{Natural density against upper density}\label{gpi:sec:gap}
{\raggedright\sources{grh_ideal_densities_and_transfer_conditions_2026-11-26; grh_transfers_density_gaps_and_zero_tests; grh_transfer_advances_2026-11-26}}

\begin{theorem}[Separation in every number field; \stU]\label{gpi:thm:ideal-separation}
For every number field $K$ and for either convention, $u_K>b_K$. More precisely, fix a prime
ideal $\mathfrak p$, put $Q=\Nm\mathfrak p$ and $M=Q^8$, and for an integer $Y\ge Q$ let
$c_Y=\prod_{\Nm\mathfrak q\le Y}(1-w_{\mathfrak q})$. Then
\begin{equation}\label{gpi:eq:ideal-gap}
 u_K-b_K\ \ge\ \frac{c_Y}{M(M+1)}-\frac nY ,
\end{equation}
and the right side is positive for every sufficiently large $Y$. Both $b_K$ and $u_K$ are
computable from an explicit presentation of $K$.
\end{theorem}

\begin{proof}
Use the notation of Lemma~\ref{gpi:lem:BU}. Let $G$ be the greedy smooth set, consisting of the
smooth ideals all of whose exponents lie in $\mathcal A$. For $m=\sum_j\epsilon_j3^j$ with
digits $\epsilon_j\in\set{0,1,2}$ define
\[
 a(m)=\sum_{\epsilon_j=1}3^j,\qquad b(m)=\sum_{\epsilon_j\in\set{1,2}}3^j .
\]
Then $a(m),b(m)\in\mathcal A$ and $m=2b(m)-a(m)$; if $m\in\mathcal A$ both equal $m$, and
otherwise $0\le a(m)<b(m)<m$. Applying this in every coordinate shows that every nongreedy
ideal is the last term of a progression with integral ratio whose first two terms are greedy
ideals of strictly smaller norm. On the other hand all greedy ideals coexist without any
forbidden triple, even with fractional ratios.

It follows by induction on the integer norm that an admissible set of smooth ideals which
maximizes the number of its elements of norm at most $k$ for \emph{every} integer cutoff
$k=1,2,\dots,M$ must agree with $G$ up to norm $M$: in a new norm layer every nongreedy
candidate is already blocked, all greedy candidates can be included together, and to reach
the greedy count all of them must be included. The argument needs no ordering inside a norm
layer, and therefore also covers distinct ideals of equal norm.

Now change only the pure powers of $\mathfrak p$ below the cutoff $M=Q^8$: replace the
exponent set $\set{0,1,3,4}$ by $\set{0,1,3,7,8}$. The latter contains no three-term
arithmetic progression and has one more element. The resulting set of smooth ideals of norm
at most $M$ is admissible even for fractional ratios. Indeed, a new forbidden triple has an
added pure power of $\mathfrak p$ as one of its terms. At every other prime one of the three
valuations is then zero and the other two lie in $\mathcal A$; since $0\in\mathcal A$ and
$\mathcal A$ is progression-free, all three are zero. The triple is confined to pure powers
of $\mathfrak p$, where the new exponent set excludes it.

Thus $G$ is not of maximal cardinality at the cutoff $M$, and no admissible set is optimal at
all cutoffs up to $M$. For every admissible smooth $D$ there is an integer $1\le k\le M$ with
$J(k)-D(k)\ge1$. The deficit persists on $[k,k+1)$, so
\[
 \int_1^\infty\bigl(J(t)-D(t)\bigr)\frac{\dd t}{t^2}\ \ge\ \frac1{k(k+1)}\ \ge\ \frac1{M(M+1)},
\]
and \eqref{gpi:eq:gapframe} gives $U-B\ge c/(M(M+1))$ and then \eqref{gpi:eq:ideal-gap}, by
the tail bound of Lemma~\ref{gpi:lem:BU}. Finally
$c_Y\ge\prod_{p\le Y}(1-p^{-1})^n\gg_n(\log Y)^{-n}$ by Mertens' product theorem, which
decreases more slowly than $n/Y$. This proves positivity.

For computability choose $Y$ with $n/Y$ below the desired tolerance; by
Lemma~\ref{gpi:lem:BU}, $0\le B-b_K\le n/Y$ and $0\le U-u_K\le n/Y$. The smooth reciprocal sum over
norms at most $H$ approaches $1/c$, with the explicit rational tail
$\tau_H=1/c-\sum_{\Nm\mathfrak h\le H}1/\Nm\mathfrak h$. A finite weighted independent-set
optimization over the smooth ideals of norm at most $H$ approximates $B$ within $c\tau_H$.
Optimizing cardinalities at each norm up to $H$ and integrating the resulting step function
approximates $U$ within $c\int_H^\infty H_S(t)t^{-2}\,\dd t=c\,(H_S(H)/H+\tau_H)$. Both errors
are computable rationals tending to zero as $H\to\infty$. (This is the existence of an
algorithm; the exact independent-set problems can be expensive.)
\end{proof}

\begin{example}[The field $\Q(\sqrt{-23})$]\label{gpi:ex:Q23gap}
Let $K=\Q(\sqrt{-23})$ and choose a prime of norm $2$, so $M=256$. A rational prime $p\ne23$
splits or is inert according to its quadratic-residue class modulo $23$; $23$ ramifies, split
primes contribute two Euler factors to $c_Y$, and inert primes contribute one factor when
$p^2\le Y$. For $Y=10^7$ there are $664449$ factors, and rounding down on the grid
$10^{-40}\Z$ after each factor gives
\[
 c_{10^7}\ \ge\ \frac{11078476589752147709803027083658225201}
 {625000000000000000000000000000000000000}.
\]
With the tail bound $2/10^7$, \eqref{gpi:eq:ideal-gap} yields
$u_K-b_K>1/20000000$ for both conventions. The computation is
\gpifile[,]{grh_ideal_densities_and_transfer_conditions_2026-11-26/code/check_ideal_density.py}
with the sharper rational bound recorded in
\gpifile[.]{grh_ideal_densities_and_transfer_conditions_2026-11-26/data/ideal_density_certificate.json}
\end{example}

For fractional ratios \cref{gpi:thm:ideal-separation} shows in particular that the greedy
density is exceeded by an upper density; that it is also exceeded by a natural density is
\cref{gpi:thm:onedigit}.

For $K=\Q$ the bound \eqref{gpi:eq:ideal-gap} is far from the truth, and the same frame gives
much more when the prefix losses are computed exactly. Let $K=\Q$, fix a convention and an
integer $N\ge1$, and let $J(k)$ now denote the maximal cardinality of an admissible subset of
$\set{1,\dots,k}$. Put
\begin{equation}\label{gpi:eq:LN}
 L_N=\sum_{k=1}^N\frac{J(k)}{k(k+1)}
 -\max_{A\subseteq\set{1,\dots,N}\ \mathrm{admissible}}\
 \sum_{a\in A}\Bigl(\frac1a-\frac1{N+1}\Bigr).
\end{equation}

\begin{lemma}[Prefix losses]\label{gpi:lem:prefix}
For $K=\Q$ and every integer $Y\ge N$ one has $U-B\ge c\,L_N$, and hence, for the respective
convention,
\[
 u_\Q-b_\Q\ \ge\ L_N\prod_{p\le Y}(1-p^{-1})-\sum_{p>Y}p^{-2}.
\]
\end{lemma}

\begin{proof}
For $Y\ge N$ all integers up to $N$ are $Y$-smooth, so the function $J$ of
\cref{gpi:sec:local} agrees on $[1,N+1)$ with the one used in \eqref{gpi:eq:LN}. For an
admissible $A\subseteq S_Y$ the second expression in \eqref{gpi:eq:LN}, evaluated at
$A\cap[1,N]$, is $\int_1^{N+1}\#(A\cap[1,t])\,t^{-2}\,\dd t$, and the first is
$\int_1^{N+1}J(t)t^{-2}\,\dd t$. The remaining integral of the prefix deficit over
$[N+1,\infty)$ is nonnegative. Now use \eqref{gpi:eq:gapframe}.
\end{proof}

\begin{example}[The obstruction in $\set{1,\dots,12}$]\label{gpi:ex:twelve}
The rational-ratio progressions up to $12$ are
\[
 (1,2,4),\quad(1,3,9),\quad(2,4,8),\quad(3,6,12),\quad(4,6,9),
\]
as one sees from the parametrization $(ab^2,abc,ac^2)$ with $1\le b<c$,
$\gcd(b,c)=1$ and $ac^2\le12$; the first four have integer ratios. An admissible set that
maximizes its cardinality at every cutoff $k\le12$ must agree with the increasing greedy
construction up to $12$: whenever the greedy construction takes $k$, adding $k$ to the agreed
prefix improves the count, so every simultaneous maximizer takes it, and whenever it rejects
$k$, taking $k$ would violate admissibility. Under both conventions the greedy prefix is
$\set{1,2,3,5,6,7,8,10,11}$, with nine elements, whereas
$\set{1,2,5,6,7,8,9,10,11,12}$ is admissible and has ten. Hence every admissible set falls
short of $J(k)$ at some $k\le12$, and $L_{12}\ge1/(12\cdot13)$ for both conventions.
\end{example}

\begin{theorem}[McNew's Questions 7.1 and 7.3; \stU]\label{gpi:thm:integer-gap}
\[
 \overline\alpha-\alpha>0.002836,\qquad\overline\beta-\beta>0.069218 .
\]
\end{theorem}

\begin{proof}
For $N=80$ the exact values are
\[
 L_{80}=\frac{1553}{45360}\ \text{(rational ratios)},\qquad
 L_{80}=\frac{16567}{132300}\ \text{(integer ratios)}.
\]
The finite maximum is computed as follows. Form each forbidden triple in $[1,k]$ as an edge
with three vertices; maximizing the retained weight is the same as minimizing the weight of a
set of vertices meeting every edge. If an uncovered edge is $\set{i,j,l}$, branch on putting
each of its three vertices into the cover and remove all incident edges; the optimum is the
minimum over the branches, with zero cost for the empty edge set, and identical remaining
edge sets are cached. The same recursion with unit weights gives each $J(k)$, and with integer
multiples of the weights in \eqref{gpi:eq:LN} it gives the second term. With $Y=293$ and
rational ratios, Lemma~\ref{gpi:lem:prefix} gives $\overline\alpha-\alpha>0.002836$; the product
is bounded below and the tail $\sum_{p>Y}p^{-2}$ is bounded above by sieving the primes up to
$100000$ and adding $1/100000$ for the rest, in directed fixed-point integer arithmetic, by
\gpifile[.]{grh_transfers_density_gaps_and_zero_tests/checks/check_integer_prefix.py}
The same method with integer ratios and $Y=83$ gives $\overline\beta-\beta>0.013336$.

For integer ratios a separate calculation is stronger: $\beta<0.749192$ by
\eqref{gpi:eq:beta-upper} in the proof of \cref{gpi:prop:beta} below, and McNew's
Theorem~4.4 gives $\overline\beta>0.818410$ \cite{McNew2015Sets}. The difference exceeds
$0.069218$.
\end{proof}

\subsection{Bounds for \texorpdfstring{$\beta$}{beta} and for the upper density
\texorpdfstring{$\overline\alpha$}{alpha-bar}}\label{gpi:sec:bounds}
{\raggedright\sources{mcnew_density_and_grh_transfers; mcnew_counterexample_and_grh_transfers; grh_transfers_density_gaps_and_zero_tests}}

\begin{proposition}[Natural density with integer ratios; \stU]\label{gpi:prop:beta}
\[
 0.723891309687823565\ \le\ \beta\ <\ 0.749192 .
\]
\end{proposition}

\begin{proof}
For the upper bound apply Lemma~\ref{gpi:lem:BU} with the integral convention and $Y=7$, so
that $c=8/35$. There are $265$ integers in $S_7\cap[1,5000]$, and the maximal reciprocal sum
of a subset of them avoiding integer-ratio progressions is
\[
 M_0=\frac{1619153731921}{497871360000}.
\]
Assigning the full remaining smooth mass to the untested tail gives
\begin{equation}\label{gpi:eq:beta-upper}
 \beta\ \le\ B\ \le\ 1-c\Bigl(\sum_{h\in S_7,\ h\le5000}\frac1h-M_0\Bigr)
 =\frac{5245322990569}{7001316000000}<0.749192 .
\end{equation}
The exhaustive calculation uses the common denominator of the $265$ weights. In ascending
order it branches on including or excluding the least undecided integer; whenever the first
two terms of a forbidden progression have been selected, the third is removed from the
candidates, and the selected and remaining weights together bound the optimum on a branch; this is
\gpifile[.]{grh_transfers_density_gaps_and_zero_tests/checks/check_integer_gap.py}

For the lower bound let $P$ be the set of the
eight primes up to $19$. The file
\gpifile{mcnew_density_and_grh_transfers/checks/integral_ratio_support.txt}
lists a set
$J_0$ of $8197$ $P$-smooth integers, the largest being $99{,}792{,}810$, which contains no
integer-ratio progression. Use Rankin's rule at the primes outside $P$. By
Lemma~\ref{gpi:lem:product} the resulting set is admissible for integer ratios and has natural
density
\[
 W(J_0)\prod_{p\le19}(1-p^{-1})
 \prod_{p>19}\Bigl[(1-p^{-2})\prod_{j\ge1}(1+p^{-3^j})\Bigr].
\]
Verifying the triples in $J_0$, retaining the factors $j=1$ for $p\le200000$ and bounding the
remaining product from below by $200000/200001$, with all fixed-point divisions rounded
downward, gives the stated bound, by
\gpifile[.]{mcnew_density_and_grh_transfers/checks/verify_integer_ratio.py}
Optimality of
the $8197$ choices is not needed.
\end{proof}

McNew's Theorem~3.3 gave $\beta>0.72195$ by modifying the greedy construction at three primes.
Upper bounds for $\alpha$ are in Corollary~\ref{gpi:cor:restricted} and Theorem~\ref{gpi:thm:beta-alpha}.

For the upper density with rational ratios we need a restricted variant. For $s\ge2$ let
$\overline\alpha_s$ be the supremum of the upper densities of the sets containing no
progression whose ratio, in lowest terms, has numerator and denominator in $S_s$. Since fewer
triples are forbidden, $\overline\alpha\le\overline\alpha_s$ for every $s$, and
$\overline\alpha_s$ is nonincreasing in $s$.

\begin{theorem}[Upper density with rational ratios; \stU]\label{gpi:thm:upper-alpha}
\[
 0.7365954<\overline\alpha<0.7716011,\qquad
 \overline\alpha_7-\overline\alpha_{11}\ \ge\ \frac{64}{139755}.
\]
\end{theorem}

McNew's bounds are $0.730027<\overline\alpha<0.772059$ \cite[Theorems~4.3 and~4.5]{McNew2015Sets};
smooth-number fibers were already used there. The additions are explicit finite templates
and a strict, certified loss when the prime $11$ is admitted to the ratios.

\begin{proof}[Proof of the lower bound]
The file
\gpifile{mcnew_counterexample_and_grh_transfers/checks/density/lower.csv}
gives
thresholds $1=t_1<\dots<t_{275}=972405$ and sets $C_i\subset S_7\cap[1,t_i]$ with $\#C_i=i$
such that $b^2/a\notin C_i$ for all $a<b$ in $C_i$ with $a\mid b^2$. Thus every $C_i$ is
admissible for rational ratios; the sets need not be nested. For $x\ge1$ let $C(x)=C_i$ for
the largest $t_i\le x$, and let $C(x)=\emptyset$ for $x<1$. The exact sum of the threshold
contributions is
\[
 I=\sum_{i=1}^{275}\frac1{t_i}=\frac{86690601810929969}{26132271943680000}.
\]
Call an integer $h$ \emph{allowed} if $\gcd(h,210)=1$, its exponents at the primes
$11\le p\le997$ lie in $\set{0,1,3,4}$, and its exponents at all larger primes lie in
$\set{0,1}$. Both exponent sets are free of nonconstant three-term arithmetic progressions.
By Lemma~\ref{gpi:lem:product} the allowed integers have natural density
\[
 c_0=\frac6{\pi^2}\prod_{p\in\set{2,3,5,7}}\frac p{p+1}\prod_{11\le p\le997}(1+p^{-3}),
\]
since the local factor $(1-p^{-1})(1+p^{-1}+p^{-3}+p^{-4})$ equals $(1-p^{-2})(1+p^{-3})$.

Take $N_0=0$ and $N_j=2^{2^j}$ for $j\ge1$. For each allowed $h\in(N_{j-1},N_j]$ insert the
integers $h\,C(N_j/h)$ into one set $A$. The part of each inserted integer coprime to $210$
is $h$. If three inserted integers formed a progression, their valuations at each prime
greater than $7$ would form an arithmetic progression in a permitted exponent set, so they
would have the same part $h$, and their $7$-smooth parts would lie in the one admissible
template assigned to $h$. This excludes progressions across stages as well as within a
stage. At $N=N_j$ the current stage alone contributes at least
\[
 \sum_{i=1}^{275}\#\set{h\le N/t_i:h\text{ allowed}}-275\,N_{j-1}
\]
elements, because the size of each template is the number of thresholds below its cutoff.
Since $N_{j-1}/N_j\to0$, division by $N_j$ gives $\overline d(A)\ge c_0I$. Replacing $\pi$ by
the upper bound $355/113$, itself certified from Machin's formula
$\pi=16\arctan(1/5)-4\arctan(1/239)$, gives the rational estimate
\[
 \frac{6\cdot113^2}{355^2}\prod_{p\in\set{2,3,5,7}}\frac p{p+1}
 \prod_{11\le p\le997}(1+p^{-3})\,I>\frac{7365954}{10^7}
\]
which is checked by
\gpifile[.]{mcnew_counterexample_and_grh_transfers/checks/density/check_density.py}
\end{proof}

\begin{proof}[Proof of the loss at eleven and of the upper bound]
Let $F_s(x)$ be the maximal size of a subset of $S_s\cap[1,x]$ free of rational-ratio
progressions, and let $g_s(N)$ be the maximal size of a subset of $[1,N]$ avoiding
progressions with $s$-smooth ratio. McNew's Theorem~4.2 gives $g_s(N)/N\to\overline\alpha_s$.
Partition $[1,N]$ into fibers according to the part $b$ coprime to $2310$, and put $x=N/b$.
When all $11$-smooth ratios are forbidden the maximum in the fiber is $F_{11}(x)$. When only
$7$-smooth ratios are forbidden it is $H(x)=\sum_{j\ge0}F_7(x/11^j)$, because each
$11$-adic layer can be chosen independently; always $H(x)\ge F_{11}(x)$. The following finite
comparison shows $H(x)\ge F_{11}(x)+1$ throughout $121\le x<165$.
\[
\small\setlength{\arraycolsep}{3pt}
\begin{array}{c|rrrrrrr}
 [a,b') &[121,125)&[125,132)&[132,135)&[135,140)&[140,154)&[154,162)&[162,165)\\
 \hline
 \max(S_{11}\cap[1,b')) &121&128&132&135&150&160&162\\
 H(a)\ \text{at least} &38&39&40&41&42&43&44\\
 F_{11}(x)\ \text{at most} &37&38&39&40&41&42&43
\end{array}
\]
For each interval, the file
\gpifile{mcnew_counterexample_and_grh_transfers/checks/density/intervals.json}
supplies three
$7$-smooth sets for the layers $j=0,1,2$; testing their members and their triples proves the
lower entry for $H(a)$. For an upper entry $u$, consider the hypergraph whose vertices are
the $11$-smooth numbers up to the listed maximum, with a hyperedge for every triple $d<e<f$
with $df=e^2$. A set of $u+1$ surviving vertices would give a cover with at most
$\#V-u-1$ deleted vertices. Seven proof trees in
\gpifile{mcnew_counterexample_and_grh_transfers/checks/density/trees/}
rule this out. At an
internal node one branches on a single vertex being deleted or kept. After discarding the
covered edges, let $k$ be the remaining deletion budget. A leaf closes if $k<0$ or if an
uncovered edge has all its vertices fixed to be kept. The remaining type of leaf carries
positive integer weights on uncovered edges, with sum $W$ and with maximal total weight $D$
incident to an undecided vertex; every cover using at most $k$ further deletions covers
weight at most $kD$, and the supplied weights satisfy $W>kD$. The verifier checks these
integer inequalities, both branches at every internal node, all listed members and the
ordering of the smooth numbers. The trees have $813$, $3295$, $1559$, $1331$, $2047$, $3139$
and $2865$ nodes.

The fibers with $\gcd(b,2310)=1$ partition $[1,N]$ for both optimization problems, and every
$b$ with $N/165<b\le N/121$ costs at least one element in the $11$-smooth problem. Dividing
$g_7(N)-g_{11}(N)$ by $N$ and passing to the limit gives
\[
 \overline\alpha_7-\overline\alpha_{11}\ \ge\ \frac{\varphi(2310)}{2310}
 \Bigl(\frac1{121}-\frac1{165}\Bigr)=\frac{64}{139755}.
\]
McNew obtains $\overline\alpha_7<0.772059$ \cite[Section~4]{McNew2015Sets}. Consequently
\[
 \overline\alpha\le\overline\alpha_{11}<\frac{772059}{10^6}-\frac{64}{139755}
 =\frac{21567021109}{27951000000}<0.7716011 . \qedhere
\]
\end{proof}

\subsection{Attainment and computability of the natural-density suprema}\label{gpi:sec:attain}
{\raggedright\sources{grh_transfers_density_gaps_and_zero_tests; grh_prime_patterns_and_branch_precision}}

\begin{theorem}[The suprema are attained; \stU]\label{gpi:thm:attain}
For each of the two conventions over $\Q$, the natural-density supremum $b\in\set{\alpha,\beta}$
is a computable real number (by \cref{gpi:thm:ideal-separation}), and there is an admissible
set of natural density exactly $b$.
Moreover $b$ is also the supremum of the upper logarithmic densities
$\limsup_{X\to\infty}(\log X)^{-1}\sum_{m\in E,\,m\le X}1/m$ over all admissible sets $E$.
\end{theorem}

\begin{proof}
\emph{Computability.} This is the case $K=\Q$ of \cref{gpi:thm:ideal-separation}, where $b_\Q$
is $\alpha$ or $\beta$ according to the convention.

\emph{A uniform estimate for the largest prime factor.} Write $D(m)$ for the largest prime
factor of $m$, with $D(1)=1$, and let $F_x(y)$ be the proportion of the positive integers
$m\le x$ with $D(m)\le y$. For every fixed positive integer $a$,
\begin{equation}\label{gpi:eq:dickman}
 \sup_{y\ge1}\abs{F_{x/a}(y)-F_x(y)}\longrightarrow0\qquad(x\to\infty).
\end{equation}
Indeed, by Dickman's theorem \cite{Dickman1930Frequency,GranvilleSmooth} the distribution
functions of the ratio $\log D(m)/\log x$ converge to the continuous law with distribution
function $t\mapsto\rho(1/t)$ on $(0,1]$, extended by zero for $t\le0$ and by one for $t\ge1$.
Pointwise convergence of distribution functions to a continuous distribution function is
uniform, by choosing a finite grid of quantiles and using monotonicity. Replacing $\log x$ by
$\log(x/a)$ rescales $t$ by a factor tending to one, with uniform error tending to zero by
continuity on a compact interval and the constant tails.

\emph{Approximating rules.} Choose cutoffs $Y_j\to\infty$ and finite admissible sets
$A_j\subset S_{Y_j}$ such that the product sets $E_j$ of Lemma~\ref{gpi:lem:product}, with smooth
part in $A_j$ and squarefree $Y_j$-rough part, have densities $e_j\to b$. Put
$\epsilon_j=2^{-j}$. Each $E_j$ is approximated, in proportion on $[1,x]$ and uniformly for
integers $x\ge1$, within $\epsilon_j$ by a finite Boolean combination $C_j$ of divisibility
conditions: membership in a finite list of exact smooth parts is a finite divisibility
condition, and ignoring the squares of primes beyond $K$ costs at most $\sum_{p>K}p^{-2}$.
Inclusion--exclusion writes the indicator of $C_j$ as a finite linear combination
$\sum_au_a\mathbf 1_{a\mid m}$. For $y\ge D(a)$ the proportion of $m\le x$ with $a\mid m$ and
$D(m)\le y$ is $\floor{x/a}F_{x/a}(y)/\floor{x}$, which by \eqref{gpi:eq:dickman} equals
$a^{-1}F_x(y)+o_a(1)$ as $x\to\infty$, uniformly in such $y$.

\emph{Pasting by the largest prime factor.} Choose increasing thresholds $m_j>Y_j$ so large
that, for all $x,y\ge m_j$, the proportion of $m\le x$ in $C_j$ with $D(m)\le y$ differs from
$d(C_j)F_x(y)$ by at most $\epsilon_j$; using $E_j$ and $e_j$ instead costs at most
$2\epsilon_j$ more. Let $E$ contain $m$ exactly when $m_j<D(m)\le m_{j+1}$ and $m\in E_j$
for some $j\ge1$; the integers with $D(m)\le m_1$ are omitted. The set $E$ is admissible. In
a forbidden triple, consider the largest prime $\ell$ dividing any of its three terms.
Wherever $\ell$ is present it is the largest prime factor of that term, exceeds the cutoff
$Y_j$ of the rule that the term follows, and therefore occurs to exponent one. The three
valuations at $\ell$ form an arithmetic progression in $\set{0,1}$ with one value equal to
one, so all equal one. Hence the three terms have the same largest prime factor, follow the
same rule $E_j$, and cannot form a forbidden triple.

For each $j$ the proportion of $m\le x$ in $E$ with $m_j<D(m)\le m_{j+1}$ differs from $e_j$
times the proportion of that bin by at most $6\epsilon_j$; both are zero if $x<m_j$. The mass
of finitely many initial bins tends to zero, because $F_x(m)\to0$ for fixed $m$. The errors
of the bins $j\ge j_0$ sum to at most $6\sum_{j\ge j_0}\epsilon_j$, and
$\sup_{j\ge j_0}\abs{e_j-b}\to0$. Letting $x\to\infty$ and then $j_0\to\infty$ shows that $E$
has natural density $b$.

\emph{Logarithmic density.} By the last assertion of Lemma~\ref{gpi:lem:BU} every admissible set
has upper logarithmic density at most $B(Y)$ for every $Y$, hence at most $b$. A set of
natural density $b$ has logarithmic density $b$ by partial summation.
\end{proof}

The polynomial analogue, with an elementary substitute for Dickman's theorem and a decidable
extremal set, is \cref{gpi:thm:poly-natural} in \cref{gpi:sec:poly}.

\subsection{The greedy set is not optimal: an exchange in one ternary digit}\label{gpi:sec:onedigit}
{\raggedright\sources{grh_prime_patterns_and_branch_precision; mcnew_counterexample_and_grh_transfers; mcnew-density-and-dirichlet-grh}}

McNew's Conjecture~7.2 asserts $\alpha=d(R)$. It fails, and it fails in the same way in every
number field and for the stronger (fractional) convention. The mechanism is to correlate the
exponents at different primes: one forbidden digit at a \emph{marker} prime is shared by
additions indexed by many \emph{satellite} primes, and one common deletion protects all of
them. The added reciprocal mass grows like a sum of reciprocals of primes, while the
deletion is paid once.

A \emph{cap} is a subset of $\F_3^{r}$ in which $x+y+z=0$ implies $x=y=z$; we use the
representatives $0,1,2$. The binary cube $\set{0,1}^{r}$ is a cap. Fix a prime ideal
$\mathfrak l$ of $K$ (the marker) and distinct further prime ideals
$\mathfrak q_1,\dots,\mathfrak q_k$ (the satellites); write $S$ for this set of $k+1$ primes,
$w=w_{\mathfrak l}$ and $u_i=w_{\mathfrak q_i}$. Coordinates in $\F_3^{k+1}$ are indexed by
$(\mathfrak l,\mathfrak q_1,\dots,\mathfrak q_k)$. Let $\mathbf b=(1,\dots,1)$ be the
all-ones satellite vector and $\mathbf b_i$ the vector obtained from it by replacing the
$i$th entry by $2$. Put $\mathcal C=\set{0,1}^{k+1}$ and
\begin{equation}\label{gpi:eq:capC}
 \mathcal C'=\bigl(\mathcal C\setminus\set{(0,\mathbf b),(1,\mathbf b)}\bigr)\cup
 \set{(2,\mathbf b),(2,\mathbf b_1),\dots,(2,\mathbf b_k)} .
\end{equation}
Let $E$ be the set of ideals $\mathfrak a$ such that the exponent vector
$(v_{\mathfrak p}(\mathfrak a))_{\mathfrak p\in S}$ has the form $3h+c$ with
$c\in\mathcal C'$ and $h\in\mathcal A^{k+1}$, and such that
$v_{\mathfrak p}(\mathfrak a)\in\mathcal A$ for every $\mathfrak p\notin S$. Replacing
$\mathcal C'$ by $\mathcal C$ gives $\mathcal G_K$; only the units digits in base three at
the primes of $S$ differ from the greedy rule.

\begin{theorem}[One-digit exchange; \stU]\label{gpi:thm:onedigit}
With the notation above, the following hold.
\begin{enumerate}[label=\textup{(\alph*)}]
\item The set $\mathcal C'$ is a cap, and $E$ is admissible for the fractional convention.
\item $E$ has a natural density, and
\begin{equation}\label{gpi:eq:capgain}
 d(E)-d(\mathcal G_K)=\Bigl(\prod_iu_i\Bigr)
 \Bigl[w^2\Bigl(1+\sum_iu_i\Bigr)-(1+w)\Bigr]
 \prod_{\mathfrak p\in S}\Bigl[(1-w_{\mathfrak p})\sum_{h\in\mathcal A}w_{\mathfrak p}^{3h}\Bigr]
 \prod_{\mathfrak p\notin S}\Bigl[(1-w_{\mathfrak p})\sum_{e\in\mathcal A}w_{\mathfrak p}^{e}\Bigr].
\end{equation}
\item For every $K$ and every marker there are finitely many satellites with
 $\sum_iu_i>w^{-2}+w^{-1}-1$, and then $d(E)>d(\mathcal G_K)$. Hence in every number field
 the natural-density supremum for fractional ratios, and a fortiori for integral ratios,
 exceeds the greedy density.
\item For $K=\Q$, take the marker $2$ and the first odd primes until their reciprocals sum to
 more than $5$; all odd primes up to $\ceil{\exp(\exp7)}$ suffice. Then $d(E)>d(R)$. In
 particular $\alpha>d(R)$, and Conjecture~7.2 of \cite{McNew2015Sets} is false.
\end{enumerate}
\end{theorem}

\begin{proof}
(a) Consider $x+y+z=0$ in $\mathcal C'$. Three old points lie in the cap $\mathcal C$. If two
of the points are new, the marker digit of the third is $2$, so it is new as well; all three
then have satellite digits in $\set{1,2}$, which is a cap in each coordinate, and equal
marker digits. Suppose exactly one point is new. At every satellite coordinate where its
digit is $1$, the two old digits lie in $\set{0,1}$ and sum to $2$ modulo three, so both are
$1$. At its exceptional coordinate, if it has one, the old digits sum to $1$, so they are $0$
and $1$. Thus at least one old point has satellite vector $\mathbf b$, and both old points
with that satellite vector were removed.

In a forbidden triple of ideals in $E$, with any integral step in each coordinate, the three
residue vectors modulo $3$ on $S$ satisfy $x+z\equiv2y$, that is $x+y+z\equiv0$. By the cap
property they are equal. Subtracting the common vector and dividing by three gives an
arithmetic progression in $\mathcal A$ in each coordinate, which is constant. The valuations
outside $S$ are constant for the same reason.

(b) Give a digit vector $c$ the mass $W(c)=w^{c_{\mathfrak l}}\prod_iu_i^{c_i}$. By
Lemma~\ref{gpi:lem:product}, applied with the greedy rule outside $S$ and with $F$ the set of
$S$-smooth ideals whose exponent vector is $3h+c$ with $c\in\mathcal C'$, respectively
$c\in\mathcal C$, and $h\in\mathcal A^{k+1}$, both densities exist and their difference is
$W(\mathcal C')-W(\mathcal C)$ times the two products in \eqref{gpi:eq:capgain}. The removed
vectors have mass $(1+w)\prod_iu_i$ and the added ones $w^2(1+\sum_iu_i)\prod_iu_i$.

(c) The sum $\sum_{\mathfrak p}w_{\mathfrak p}$ diverges: the Dedekind zeta function has a
pole at $1$, and in the logarithm of its Euler product the prime powers of exponent at least
two contribute a bounded amount as $s\downarrow1$. The two products in \eqref{gpi:eq:capgain}
are positive.

(d) With $w=1/2$ the bracket is $(\sum_iu_i-5)/4$. Euler's divergence theorem ensures that
the initial segment is finite. For the explicit cutoff $Y=\ceil{\exp(\exp7)}$, the finite Euler
product over $p\le Y$ exceeds $\sum_{m\le Y}1/m\ge\log Y$, and in its logarithm the prime
powers of exponent at least two contribute at most $\sum_{m\ge2}1/(m(m-1))=1$. Subtracting
the contribution $1/2$ of the prime $2$ leaves at least $7-1-\tfrac12>5$.
\end{proof}

Part~(c) gives an affirmative answer, in every number field and for both conventions, to the
natural-density form of the questions of Best et al.\ recalled in \cref{gpi:sec:setting}.
The same exchange over $\F_q[x]$ is \cref{gpi:prop:poly-greedy} in \cref{gpi:sec:poly}.

\paragraph{Which exchanges with two deletions can help.}
The exchange \eqref{gpi:eq:capC} deletes two vertices of the binary cube. This is the least
possible number, and the profitable two-deletion exchanges can be classified exactly. Let
$\mathcal C=\set{0,1}^r\subset\F_3^r$ carry the weights $W(x)=\prod_iw_i^{x_i}$ with
$0<w_i\le1/2$, extended to $\F_3^r$ by the representatives $0,1,2$. Require that a
replacement cap keep all cube vertices except those of a specified set $D$ and omit every
vertex of $D$.

\begin{theorem}[Two deletions; \stU]\label{gpi:thm:edge}
If $D$ is empty, or a single vertex, or two vertices at Hamming distance at least two, then
the weight that can be added outside the cube is at most the deleted weight $W(D)$; there is
no net gain. If $D$ consists of two adjacent vertices, call their varying coordinate the
marker, of weight $w$, and their common remaining vector $b$; for a satellite coordinate $i$
put $s_i=w_i^{2-b_i}$. Then the maximal weight that can be added is
\[
 W(b)\max\Bigl\{(1+w)\max_is_i,\ w^2\Bigl(1+\sum_is_i\Bigr)\Bigr\},\qquad
 W(b)=\prod_{i\ne\mathrm{marker}}w_i^{b_i},
\]
with zero for an empty maximum, and the deleted weight is $(1+w)W(b)$.
\end{theorem}

\begin{proof}
An outside point with $h$ entries equal to $2$ lies on exactly $2^{h-1}$ lines through pairs
of cube points: at a coordinate with entry $2$ the two binary digits are opposite, and
everywhere else both equal the digit of the outside point. These pairs are disjoint. Hence
$h\le2$ after two deletions. For $h=2$ both deleted vertices must belong to the square face
of the point and must meet the two different opposite pairs; in a binary square this means
that they are adjacent. For $h=1$ at least one endpoint of the associated cube edge must
have been deleted. With no deletion every outside point is excluded.

For a single missing vertex $v$ the allowed additions therefore arise by changing one digit
of $v$ to $2$. Two different such changes have as third point the cube vertex obtained from
$v$ by flipping both binary digits, which survives. So at most one point can be added, of
weight at most $W(v)/2$. For two nonadjacent missing vertices each possible addition is
incident to exactly one of them. Two different additions incident to $v$ require that their
two-digit flip equal the other missing vertex, so at most two additions are incident to $v$,
each of weight at most $W(v)/2$. Summing the charges proves the first claim.

Now let $D$ be an edge. The available points are the centre $(2,b)$ and, for each satellite
$i$, the three points with any marker digit above the satellite vector $b_i$ obtained from
$b$ by changing its $i$th digit to $2$; call this triple column $i$. The face argument gives
this entire list and shows that each individual candidate is compatible with all pairs of
retained cube points.

Suppose different columns $i,j$ are used with marker digits $a,c$. The third point of the
line has satellite vector $b$ flipped at $i$ and $j$, which is binary; its marker $-a-c$
must be $2$, since otherwise it is a retained cube vertex. Thus $a+c=1$ modulo three. A
centre and a column point give a one-flip binary completion by the same calculation, which
forces the column marker to be $2$. Three or more occupied columns force every marker to be
$2$. Two occupied columns without the centre can only carry the marker pairs $(0,1)$, $(1,0)$,
$(2,2)$, with a single point in each column. One occupied column contains at most two
points, because its three points form a line; their combined relative marker weight is at
most $1+w$. The two mixed-marker cases cost at most $(1+w)\max_is_i$ in units of $W(b)$ as
well, and all cases with marker $2$ only cost at most the second term of the maximum.

Both bounds are attained when the corresponding candidates exist. In a maximizing column use
the markers $0$ and $1$; their completion is the third point of the same column and is
absent. Alternatively use the centre and all columns with marker $2$: the satellite digits
then lie in the cap $\prod_i\set{b_i,2}$, and with only one new point a pair of retained
points cannot complete a line. This proves the formula, including the case of no satellites.
\end{proof}

The proof is confirmed on $378$ small weighted cube instances by
\gpifile[.]{grh_prime_patterns_and_branch_precision/checks/check_caps.py}
Three consequences
follow. First, since $u_i\le1/2$, the first term of the maximum never exceeds the deleted
weight, so a strict gain requires
\begin{equation}\label{gpi:eq:edge-threshold}
 w^2\Bigl(1+\sum_is_i\Bigr)>1+w ,
\end{equation}
and by \cref{gpi:thm:edge} two is the smallest number of deleted cube vertices in a profitable
one-digit exchange. The restriction concerns the residue cap and allows arbitrarily many
added patterns. Second, since $s_i\le w_i$ with equality for $b_i=1$, for a fixed marker the
most favourable choice is the all-ones satellite vector of \eqref{gpi:eq:capC}. Third, the
marker giving the easiest threshold is one of largest weight: with $S_0=\sum_iw_i$ fixed over
all coordinates, the ratio $w^2(1+S_0-w)/(1+w)$ increases for $0<w\le1/2$, since its
logarithmic derivative is $2/w-1/(1+S_0-w)-1/(1+w)>0$. Consequently the threshold
$\sum_iu_i>5$ of \cref{gpi:thm:onedigit}(d) is sharp for a two-deletion exchange with marker
$2$, and for a list of rational primes containing $2$ no other marker lowers it.

\paragraph{The least number of primes as a function of the field degree.}
Order the multiset containing $n$ copies of $1/p$ for each rational prime $p$ decreasingly,
and let $\Sigma_n(r)$ be the sum of its first $r$ terms.

\begin{proposition}[Sharp dimension; \stU]\label{gpi:prop:thresholds}
Among all number fields of degree $n$, the least number of prime-ideal coordinates in which
a two-deletion exchange strictly improves the greedy cube is
\[
 D_n=\min\set{r:\Sigma_n(r)>11/2}.
\]
Every field of degree $n$ needs at least $D_n$ coordinates, and at least one field of degree
$n$ attains $D_n$. The values are
\begin{center}
\begin{tabular}{c|rrrrrrrrrr}
$n$&2&3&4&5&6&7&8&9&10&$\ge11$\\ \hline
$D_n$&30\,948&85&27&18&15&14&13&13&12&12
\end{tabular}
\end{center}
and for $n=1$
\[
 3.46\cdot10^{79}<D_1<3.52\cdot10^{79}.
\]
\end{proposition}

\begin{proof}
By \eqref{gpi:eq:edge-threshold} and $s_i\le w_i$, the sum of all the prime-ideal weights in
the exchange, marker included, must exceed $w+(1+w)/w^2-1$, a decreasing function of $w$ on
$(0,1/2]$ with minimum $11/2$. There are at most $n$ primes above each rational prime $p$,
each of norm at least $p$, so $\Sigma_n(r)$ dominates the sum of any $r$ weights in a field of
degree $n$. The minimum is finite by the divergence of $\sum_p1/p$.

For the converse, choose a field of degree $n$ in which every rational prime occurring among
the first $D_n$ terms of the multiset splits completely. Such a field exists by coefficientwise
Chinese remaindering: set $g(X)=\prod_{j=1}^n(X-j)$ and, at every specified prime $p$, demand
that a monic polynomial $f$ of degree $n$ be congruent to $g$ modulo $p^{2M+1}$ with
$M=\max_jv_p(g'(j))$. Hensel's lemma at each $j$ gives a root whose difference from $j$ has
valuation greater than $M$, and these roots are distinct because $v_p(j-k)\le M$ for $j\ne k$.
Impose Eisenstein's criterion at one further prime to guarantee irreducibility. Then
$\Q[X]/(f)$ has the required degree and splitting, as one sees from its tensor products with
$\Q_p$. In these coordinates choose a marker of norm two and all satellite digits equal to
one; the sum condition is sufficient.

For $n\ge12$ the first twelve terms are all $1/2$; for $n=2$ the last prime used is $169751$.
The table is reproduced, with integer interval certificates and exact short sums, by
\gpifile[.]{grh_prime_patterns_and_branch_precision/checks/check_field_thresholds.py}

For $n=1$ put $x_-=6.51\cdot10^{81}$ and $x_+=6.57\cdot10^{81}$. By Theorem~5 of Rosser and
Schoenfeld \cite{RosserSchoenfeld1962Approximate}, the absolute error in replacing $\sum_{p\le x}1/p$ by
$\log\log x+B_1$ is at most $1/(2\log^2x)$ at both points, where the Meissel--Mertens
constant satisfies $0.261496<B_1<0.261499$. Substitution places the two prime sums on opposite
sides of $11/2$. Theorem~1 of the same paper bounds $\pi(x)$ below and above by
$(x/\log x)(1+\theta/\log x)$ with $\theta=1/2$ and $3/2$ respectively in this range. Applying
these bounds to $\pi(x_-)<D_1\le\pi(x_+)$ gives the interval. The checker also bounds $B_1$
directly by truncating $B_1=\gamma-\sum_p\sum_{k\ge2}1/(kp^k)$ with explicit tail errors.
\end{proof}

\begin{example}[A degree-ten field with twelve primes]\label{gpi:ex:431}
Let $K$ be the unique subfield of degree $10$ of $\Q(\zeta_{431})$. The integer $431$ is
prime, the Galois group of $\Q(\zeta_{431})$ is cyclic of order $430$, and the subgroup fixing
$K$ has order $43$. Repeated squaring gives $2^{43}\equiv3^{43}\equiv1\pmod{431}$, so $2$ and
$3$ split completely in $K$. Take the ten primes of norm $2$ and two primes of norm $3$, with a
marker of norm $2$. The eleven satellite weights sum to $9/2+2/3=31/6$, and
\eqref{gpi:eq:capgain} gives an added unnormalized mass of $1/110592$. Dividing by the mass
$(3/2)^{10}(4/3)^2$ of the original cube gives the exact relative gain
\[
 d(E)=d(\mathcal G_K)\Bigl(1+\frac1{11337408}\Bigr).
\]
By Proposition~\ref{gpi:prop:thresholds}, twelve is the least possible number of coordinates for a
two-deletion exchange in any number field, and ten is the least degree in which twelve
suffice. The arithmetic of this example is checked by
\gpifile[.]{grh_prime_patterns_and_branch_precision/checks/check_fields.py}
\end{example}

\begin{example}[No uniform bound for the relative gain]\label{gpi:ex:unbounded}
There is no field-independent upper bound for the ratio of the achievable natural density to
the greedy density. Choose $L$ distinct rational primes congruent to $1$ modulo $8$ and
adjoin their square roots. Independent square classes give a field of degree $2^L$, and all
radicands are squares in $\Q_2$, so $2$ splits completely. Apply the exchange
\eqref{gpi:eq:capC} in disjoint groups of twelve prime ideals of norm $2$. Each group
multiplies the density by exactly $1+1/2125764$, since its added mass is $1/16384$ and its
original mass is $(3/2)^{12}$, and the admissibility proof applies to each group separately.
Independent groups multiply the gains, giving the factor
$(1+1/2125764)^{\floor{2^L/12}}$ relative to the greedy density of the same field.
\end{example}

\subsection{An explicit counterexample over the integers}\label{gpi:sec:explicit}
{\raggedright\sources{mcnew_counterexample_and_grh_transfers; mcnew-density-and-dirichlet-grh}}

By Proposition~\ref{gpi:prop:thresholds} the one-digit exchange over $\Q$ needs about $3.5\cdot10^{79}$
primes. An exchange of exact valuations, with four marker exponents, needs far fewer.

\begin{theorem}[A counterexample on the primes up to $1427$; \stU]\label{gpi:thm:explicit}
There is a set $E\subseteq\N$ containing no three distinct integers $a,b,c$ with $ac=b^2$,
having a natural density, which differs from Rankin's set only in the exponents at the $225$
primes up to $1427$, and such that
\[
 d(E)-d(R)>10^{-596}.
\]
\end{theorem}

\begin{proof}
Let $B$ be the set of the $224$ odd primes $q\le Y=1427$, put $Q=\prod_{q\in B}q$ and
$S=\set{2}\cup B$, and let $C$ be the set of $S$-smooth integers all of whose exponents lie
in $\mathcal A$. For an $S$-smooth integer $m$ write $a=v_2(m)$ and $b_q=v_q(m)$. Let
\[
 \mathcal K=2+3\set{0,1,3,4}=\set{2,5,11,14},\qquad
 2\mathcal K=\set{2k:k\in\mathcal K}.
\]
Delete $m\in C$ under either of the following conditions.
\begin{itemize}
\item All odd exponents $b_q$ equal $1$ and $a>0$.
\item Exactly one odd exponent $b_q$ differs from $1$, and either $b_q=4$, or $b_q=3$ and
 $a\in2\mathcal K$.
\end{itemize}
Insert the integers $2^kQ$, $2^kQq$ and $2^kQq^3$ for every $q\in B$ and $k\in\mathcal K$.
All four markers $k$ are congruent to $2$ modulo $3$, so the inserted integers lie outside
$C$. Let $F$ be the resulting set of $S$-smooth integers and let $E$ be the product set of
Lemma~\ref{gpi:lem:product} with smooth part in $F$ and Rankin's rule outside $S$.

\emph{A local fact.} If $x,y\in\mathcal A$ and $x,2,y$ is an arithmetic progression, then
$\set{x,y}=\set{0,4}$ or $\set{1,3}$. If $x,y,2$ is an arithmetic progression with
$x,y\in\mathcal A$, then $x=2y-2$ and $y\in\set{1,3}$. For the latter write $y$ in base
three. If its last digit is $1$, subtracting $2$ from $2y$ leaves every higher digit doubled,
so all higher digits vanish. If its last digit is $0$, let $j\ge1$ be the position of its
first nonzero digit; borrowing to subtract $2$ would create a digit $2$ at position $1$ if
$j>1$, so $j=1$, and all higher digits vanish, again by doubling. Both values of $y$ occur.

\emph{Admissibility of $F$.} Three points of $C$ form no forbidden triple. Consider a
forbidden triple containing exactly one inserted point.

If it is $2^kQq$, then at every $q'\in B\setminus\set q$ the three exponents lie in
$\mathcal A$ and one of them is $1$, so all are $1$. If the inserted point is an end of the
progression, the middle exponent at $q$ is $1$ or $3$ by the local fact. In the first case
the middle point has all odd exponents $1$ and a positive exponent at $2$, since that
exponent is the average of the positive $k$ and a nonnegative integer; it was deleted. In
the second case the other end has exponent $4$ at $q$ and was deleted. If the inserted point
is the middle, the outer exponents at $q$ are $0,4$ or $1,3$. A point with exponent $4$ was
deleted. In the other pair the point with exponent $1$ at $q$ has all odd exponents $1$ and
survives only for $a=0$, which forces the point with exponent $3$ at $q$ to have
$a=2k\in2\mathcal K$, so that it was deleted.

If it is $2^kQ$, both old points have all odd exponents $1$ and survive only with $a=0$;
the exponents $0,0,k$ at $2$ are not in arithmetic progression in any order.

If it is $2^kQq^3$, the exceptional odd exponent is $4\in\mathcal A$. Every odd coordinate
of all three points then lies in $\mathcal A$ and is constant, so both old points have the
deleted odd pattern with exceptional exponent $4$.

If two points are inserted, their exponents at $2$ are congruent to $2$ modulo $3$, which
forces the third exponent at $2$ into the same class in each of the three positions; no
member of $\mathcal A$ lies in it. With three inserted points every odd exponent lies in the
progression-free set $\set{1,2,4}$ and the exponent at $2$ in the progression-free set
$\mathcal K$, an affine image of a subset of $\mathcal A$. Each coordinate is then constant
and the three terms coincide. The cases allow any order of the terms, hence all rational
ratios. By Lemma~\ref{gpi:lem:product}, $E$ is admissible and has a natural density.

\emph{The gain.} Let $\sigma=\sum_{a\in\mathcal A}2^{-a}$. Exponents below $27$ use only the
three lowest ternary digits, and the remaining terms are bounded by
$\sum_{a\ge27}2^{-a}=2^{-26}$, so
\[
 \sigma\le\Upsilon:=(1+2^{-1})(1+2^{-3})(1+2^{-9})+2^{-26}.
\]
Put $u=\sum_{k\in\mathcal K}2^{-k}$ and $v=\sum_{k\in\mathcal K}2^{-2k}$. In units of $1/Q$
the insertions contribute $u\,(1+\sum_{q\in B}(q^{-1}+q^{-3}))$ to the reciprocal mass, the
common deletion costs $\sigma-1$, and the two kinds of single-coordinate deletions cost at
most $v/q^2$ and exactly $\sigma/q^3$ respectively (the first bound also counts those $2k$
that are not in $\mathcal A$). Hence
\[
 Q\,\bigl(W(F)-W(C)\bigr)\ \ge\ u-(\Upsilon-1)
 +\sum_{q\in B}\Bigl(\frac uq-\frac v{q^2}-\frac{\Upsilon-u}{q^3}\Bigr)>\frac9{100000}.
\]
The last inequality is checked in exact rational arithmetic by
\gpifile[,]{mcnew_counterexample_and_grh_transfers/checks/check_rankin_finite.py}
which also
sieves the primes, verifies $Q<3\cdot10^{590}$, and checks
\[
 \frac Y{Y+1}\prod_{p\le Y}(1-p^{-1})>\frac1{15}.
\]
By Lemma~\ref{gpi:lem:product}, $d(E)-d(R)=c_H(W(F)-W(C))$, where $c_H$ is the density of the
permitted rough parts. Its local factors at $p>Y$ are at least $1-p^{-2}$, with product at
least $\prod_{m>Y}(1-m^{-2})=Y/(Y+1)$, so $c_H>1/15$ and
\[
 d(E)-d(R)>\frac9{100000\cdot15\cdot3\cdot10^{590}}>10^{-596}. \qedhere
\]
\end{proof}

\subsection{Progressions of every length}\label{gpi:sec:klength}
{\raggedright\sources{prime_observations_and_progression_free_sets}}

The shared-deletion mechanism works for geometric progressions of every length. In this
subsection and in \cref{gpi:sec:restricted,gpi:sec:blocks}, $K=\Q$. For an integer $k\ge3$ a
\emph{$k$-progression} consists of distinct positive integers $a,ar,\dots,ar^{k-1}$ with
$r\in\Q_{>0}$, $r\ne1$; for $k=3$ these are the triples forbidden by the fractional convention.
Let $G_k\subseteq\Z_{\ge0}$ be the set obtained by scanning the nonnegative integers in order and
retaining an exponent unless it completes a nonconstant $k$-term arithmetic progression of
retained exponents, and put
\[
 \mathcal G_k=\set{n\ge1:v_p(n)\in G_k\text{ for every prime }p}.
\]
Since $G_k$ is the greedy $k$-term-progression-free set, $G_3=\mathcal A$ and $\mathcal G_3=R$; in
general $\set{0,1,\dots,k-2}\subseteq G_k$ and $k-1\notin G_k$. Valuations turn $k$-progressions
into nonconstant $k$-term arithmetic progressions of exponent vectors, so $\mathcal G_k$ is free of
$k$-progressions. It is also the greedy such set of positive integers. Indeed, if a candidate $s$
has exponents outside $G_k$, each such exponent $e$ was rejected at the end of a progression
$e-(k-1)\delta,\dots,e-\delta,e$ with $\delta>0$ and earlier entries in $G_k$; using one such
progression at every bad prime simultaneously, and keeping the good exponents constant, gives $k-1$
proper divisors of $s$ in $\mathcal G_k$ completing a $k$-progression that ends at $s$. Call these
the \emph{greedy witnesses} of $s$; induction proves the greedy characterization. By
\cref{gpi:lem:product}, $\mathcal G_k$ has natural density
$\prod_p(1-p^{-1})H_p$ with $H_p=\sum_{e\in G_k}p^{-e}$.

Brown and Gordon proved that $\mathcal G_k$ has the largest density among the sets defined by one
common progression-free exponent set at every prime, and after their Theorem~3 asked whether a set
outside this class can have larger density \cite{BrownGordon1996Sequences}; Riddell had improved on
Rankin's original benchmark for $k\ge4$ \cite{Riddell1969Sets}. For $k=3$ the answer is
\cref{gpi:thm:onedigit}. The following theorem gives it for every length.

\begin{theorem}[Beating the greedy set at every length; \stU]\label{gpi:thm:klength}
For every $k\ge3$ there is a set $A_k\subseteq\N$ free of $k$-progressions with
$d(A_k)>d(\mathcal G_k)$. It differs from $\mathcal G_k$ only in the joint rule on finitely many
primes, which can be chosen outside any prescribed finite set.
\end{theorem}

Fix $k$ and put $b=k-2$, so that $0,1,\dots,b\in G_k$ and $b+1\notin G_k$. Choose a \emph{marker}
prime $\ell$ and a finite nonempty set $B$ of other primes, the \emph{satellites}; put
$S=\set\ell\cup B$ and $C_S=G_k^{S}$, the old exponent vectors. Define $F_S$ from $C_S$ as follows.
\begin{enumerate}[label=\textup{(\arabic*)}]
\item Delete every old vector with $v_q=b$ for all $q\in B$, whatever its marker exponent in $G_k$.
\item For each $q\in B$, delete the old vectors with $v_q\ge b+2$ and every other satellite exponent
 equal to $b$.
\item For each $q\in B$, add the vector with $v_\ell=v_q=b+1$ and every other satellite exponent
 equal to $b$.
\end{enumerate}

\begin{lemma}[One shared deletion]\label{gpi:lem:kavoid}
$F_S$ contains no nonconstant $k$-term arithmetic progression of exponent vectors, for arbitrary
integral step vectors.
\end{lemma}
\begin{proof}
A progression of old vectors is constant in every coordinate. If at least two terms are new, the
marker takes the value $b+1$ twice and is constant, so every term is new; each satellite then takes
at most two values on the new vectors and is constant since $k\ge3$. Suppose exactly one term is
new, indexed by $q$. At each satellite other than $q$ all $k$ entries lie in $G_k$, including the new
entry $b$, so these coordinates are constant at $b$. In this slice the surviving old entries at $q$
lie in $[0,b-1]$: $b$ was deleted by the first rule, $b+1$ was never present, and the second rule
removed the larger entries. At $q$ the new entry is $b+1=k-1$, and a nonconstant $k$-term
arithmetic progression in $[0,k-1]$ uses all $k$ integers, including the missing value $b$; a
constant one cannot contain both an old and a new entry.
\end{proof}

Write $W=\prod_{q\in B}q^{-b}$ and $W_S(F)=\sum_{e\in F}\prod_{p\in S}p^{-e_p}$, the reciprocal mass of
the $S$-smooth integers with exponent vectors in $F$. The deletion families are disjoint, and
direct summation gives
\begin{equation}\label{gpi:eq:kgain}
 \frac{W_S(F_S)-W_S(C_S)}W=\ell^{-(k-1)}\sum_{q\in B}\frac1q
 -H_\ell\Bigl(1+\sum_{q\in B}\sum_{\substack{e\in G_k\\e\ge b+2}}q^{b-e}\Bigr).
\end{equation}
The subtracted term is at most
$(1-1/\ell)^{-1}\bigl(1+\sum_{q\in B}1/(q(q-1))\bigr)\le4$, since
$\sum_{j\ge2}1/(j(j-1))=1$. By Euler's divergence theorem $B$ can be chosen outside any prescribed
finite set with $\sum_{q\in B}1/q>5\ell^{k-1}$, and then \eqref{gpi:eq:kgain} exceeds $1$. For
example, at length $4$ with marker $2$, the satellite baseline is $2$, the new exceptional exponent
is $3$, and a reciprocal sum above $40$ suffices. This is an existence statement; the improvement
can be extremely small.

\begin{proof}[Proof of \cref{gpi:thm:klength}]
Use $F_S$ on $S$ and $G_k$ separately at every other prime. In a $k$-progression in this set every
coordinate would be constant, by \cref{gpi:lem:kavoid} and the definition of $G_k$, so the ratio
would be $1$. By \cref{gpi:lem:product}, applied with $A_p=G_k\ni0,1$ outside $S$, the set has a
natural density, and its excess over $d(\mathcal G_k)$ is
$c_S\prod_{p\notin S}(1-1/p)H_p\cdot(W_S(F_S)-W_S(C_S))>0$; each factor $(1-1/p)H_p$ is at least
$1-p^{-(k-1)}$.
\end{proof}

For $k=3$ this is a second disproof of Conjecture~7.2 of \cite{McNew2015Sets}, by a deletion that
involves all marker exponents rather than one ternary digit; the constants of
\cref{gpi:thm:onedigit,gpi:thm:explicit} are not reproduced by it. The edited rule is checked for
$3\le k\le9$ on all directed $k$-term progressions in the exponent box $[0,2k+1]^3$ at the primes
$2,3,5$, together with the truncated identity \eqref{gpi:eq:kgain} in exact rational arithmetic, by
\gpifile[.]{prime_observations_and_progression_free_sets/progressions/check_progressions.py}
Positivity for a sufficiently large prime set rests on the argument above.

\paragraph{Independent gains in residue classes.}
Write $[a]_m=\set{n\ge1:n\equiv a\pmod m}$; densities are relative to all positive integers.

\begin{theorem}[Prescribed gains in reduced classes; \stU]\label{gpi:thm:kresidues}
For $k\ge3$ and $m\ge2$ there is $\eta=\eta(k,m)>0$ such that for any real numbers
$0\le t_a\le\eta$ indexed by the reduced residue classes modulo $m$ there is a set $A$ free of
$k$-progressions with
\[
 d(A\cap[a]_m)=d(\mathcal G_k\cap[a]_m)+t_a\qquad((a,m)=1).
\]
One may require $A$ to agree with $\mathcal G_k$ on all nonreduced classes and on every class with
$t_a=0$.
\end{theorem}
\begin{proof}
Choose the marker and all satellites congruent to $1$ modulo $m$. Character orthogonality and the
Euler products give $\sum_{p\equiv1\bmod m}p^{-s}=\varphi(m)^{-1}\log\frac1{s-1}+O_m(1)$ as
$s\downarrow1$, since the nonprincipal Dirichlet $L$-functions do not vanish at $1$ and prime powers
of exponent at least two contribute boundedly \cite{Sutherland2021Dirichlet}. The reciprocal sum of
these primes diverges, so one can achieve $\Delta=W_S(F_S)-W_S(C_S)>0$.

Factor $n=hs$ with $s$ its $S$-smooth part and $h$ coprime to $S$; every such $s$ is $1$ modulo $m$.
Require the valuations of $h$ to lie in $G_k$. For reduced $a$ the set $H_a$ of eligible $h$ in
$[a]_m$ has density
\[
 \gamma=\frac{c_S}m\prod_{p\notin S,\ p\nmid m}(1-1/p)H_p>0,
\]
independent of $a$: the congruence fixes valuation zero at the primes dividing $m$, and the other
finite valuation conditions are independent of it by the Chinese remainder theorem, with the tail
estimate of \cref{gpi:lem:product} justifying the limit. Put $\eta=\gamma\Delta$ and
$\lambda_a=t_a/\eta$. In the ordered sequence $h_1<h_2<\cdots$ of $H_a$ mark $h_j$ exactly when
$\floor{j\lambda_a}>\floor{(j-1)\lambda_a}$; the marked set has density $\gamma\lambda_a$. At a marked $h$
use $F_S$ for $s$, at every other eligible $h$ use $C_S$; no rough part in a nonreduced class is
marked. In a $k$-progression the valuations outside $S$ are constant, so $h$ is fixed, and both
rules on its fibre are progression-free. Since $s\equiv1\pmod m$, the fibre lies in the class of $h$.
Scaling a set of rough parts by $s$ divides its density by $s$, the fibres of distinct smooth parts
are disjoint, and the smooth parts $s>T$ contribute at most $\sum_{s>T}1/s\to0$ in upper density.
The gain in $[a]_m$ is therefore $\gamma\lambda_a\Delta=t_a$, with the stated agreement elsewhere.
\end{proof}

\subsection{Where Rankin's rule is optimal}\label{gpi:sec:restricted}
{\raggedright\sources{grh_prime_patterns_and_branch_precision; mcnew_density_and_grh_transfers; prime_observations_and_progression_free_sets}}

The counterexamples couple many primes. For few primes Rankin's rule cannot be beaten. For a
finite set $P$ of rational primes write $S_P$ for the $P$-smooth positive integers,
$c_P=\prod_{p\in P}(1-p^{-1})$, and $C_P$ for the set of $m\in S_P$ with all exponents in
$\mathcal A$. In this subsection and the next, \emph{progression-free} means admissible for
rational ratios, that is, free of $3$-progressions, unless a length $k$ is specified. Say that \emph{Rankin's rule is optimal on $P$} if $W(F)\le W(C_P)$ for
every progression-free $F\subseteq S_P$.

\begin{proposition}[Light blocks; \stU]\label{gpi:prop:light}
Let $k\ge3$, let $1/c_P\le2$, and let $G_k^P$ be the set of $m\in S_P$ with all exponents in $G_k$.
Every $F\subseteq S_P$ free of $k$-progressions satisfies $W(F)\le W(G_k^P)$. In particular Rankin's
rule is optimal on $P$.
\end{proposition}

For $k=3$ this is the observation following Theorem~6.1 of \cite{McNew2015Sets}.

\begin{proof}
For $s\in F\setminus G_k^P$ the greedy witnesses of \cref{gpi:sec:klength}, formed within $P$, are $k-1$
proper divisors of $s$ in $G_k^P$ completing a $k$-progression that ends at $s$, so at least one of
them is missing from $F$; charge $1/s$ to one such divisor $t$. Distinct integers charged to the
same $t$ have distinct nontrivial $P$-smooth quotients $s/t$, so the charge at $t$ is at most
$(1/c_P-1)/t\le1/t$. Summing over the missing elements of $G_k^P$ gives
$W(F\setminus G_k^P)\le W(G_k^P\setminus F)$.
\end{proof}

\begin{theorem}[Blocks of at most four primes; \stU]\label{gpi:thm:fourprimes}
Rankin's rule is optimal on every set of at most four primes.
\end{theorem}

\begin{proof}
It suffices to treat four primes $q_1<q_2<q_3<q_4$: a smaller set can be padded with further
primes carrying independent Rankin rules, and the resulting inequality divided by their
positive reciprocal masses. If $q_1\ge5$, then
$\prod_iq_i/(q_i-1)\le(5/4)(7/6)(11/10)(13/12)<2$ and Proposition~\ref{gpi:prop:light} applies.
Otherwise select a base tuple $b$ from the following table; the first $f$ entries of the
tuple are fixed, $q_i=b_i$, and the others satisfy $q_i\ge b_i$.
\begin{center}\small
\begin{tabular}{lrrrr}
\toprule
Condition & Base tuple & $f$ & Prefixes & Search nodes\\ \midrule
$q_1=2,\ q_2=3$ & $(2,3,5,7)$ & 2 & 9 & 5{,}537{,}385\\
$q_1=2,\ q_2\ge5$ & $(2,5,7,11)$ & 1 & 81 & 221{,}071\\
$q_1=3$ & $(3,5,7,11)$ & 1 & 81 & 29{,}307\\ \bottomrule
\end{tabular}
\end{center}
Use the box $\mathcal D=\prod_i\set{0,\dots,m_i-1}$ of exponent vectors with $m=(9,9,3,3)$,
corresponding to the low ternary digit depths $(2,2,1,1)$. Write $\mathcal A^4$ for the good
vectors, $W_b(I)=\sum_{a\in I}\prod_ib_i^{-a_i}$ and $\lambda=21/20$. For every rectangular
prefix $\mathcal D'$ of $\mathcal D$ in which the first $f$ axes are not truncated, the
program
\gpifile{grh_prime_patterns_and_branch_precision/checks/check_four_prime_blocks.py}
certifies that every progression-free $I\subseteq\mathcal D'$ satisfies
\begin{equation}\label{gpi:eq:lowbox}
 \lambda\,W_b(I\setminus\mathcal A^4)\ \le\ W_b\bigl((\mathcal A^4\cap\mathcal D')\setminus I\bigr).
\end{equation}
The test orders the box by the integer $\prod_ib_i^{a_i}$ and branches on including or
omitting the least undecided value. Including a value $y$, when a smaller $x$ has already
been chosen, forbids $y^2/x$ whenever it lies in the box. The upper bound for the objective
on a branch is the summed weight of all chosen and undecided values, with the factor
$\lambda$ on the bad ones, and a branch is complete as soon as this bound is at most
$W_b(\mathcal A^4\cap\mathcal D')$. The objective is scaled to integers. Every
progression-free subset follows one branch, so a successful traversal certifies the
inequality.

\emph{From the base tuple to all tuples.} For $t_i=b_i/q_i\le1$ the function
$a_i\mapsto t_i^{a_i}$ on $0\le a_i<m_i$ is a nonnegative linear combination of indicators of
prefixes: the coefficient of $\mathbf 1_{a_i\le j}$ is $t_i^j-t_i^{j+1}$ for $j<m_i-1$ and
$t_i^{m_i-1}$ for the full range. A fixed coordinate uses only the full range. Multiply the
four expansions, apply \eqref{gpi:eq:lowbox} to $I\cap\mathcal D'$ on each prefix and sum.
This gives \eqref{gpi:eq:lowbox} on the full box with every $b_i$ replaced by $q_i$. This
change of quantifier matters, because a partition of the primes can pair even the smallest
prime with an arbitrarily large one.

\emph{From the box to all exponents.} Let $F\subseteq S_{\set{q_1,\dots,q_4}}$ be
progression-free and put $M=W_q(C\setminus F)$, where $C$ is the Rankin set. Split exponents
into low digits and a high part divisible by $m_i$ in each coordinate. For each fixed good
high part, \eqref{gpi:eq:lowbox} bounds the mass of the points of $F$ with a bad low digit by
the missing good mass divided by $\lambda$; dilation by a fixed high part preserves both the
weights up to a common factor and the forbidden triples. Summing gives at most $M/\lambda$.
For a point $s$ of $F$ with a bad high digit, replace every digit $2$ by $0$, respectively
by $1$; the two resulting points $a<b<s$ of $C$ satisfy $as=b^2$, so one of them is missing from
$F$. Relative to a missing point the charged reciprocal weight is a product of
factors $q_i^{-e3^j}$ over a set of digit positions, with $e=1$ for the replacement by $1$
and $e=2$ for the replacement by $0$, and at least one changed digit is high. The total
charge is at most $M\,T(q)$, where
\[
 T(q)=\sum_{e=1,2}\Bigl[\prod_i\prod_{0\le j<k_i}(1+q_i^{-e3^j})\Bigr]
 \Bigl[\prod_i\bigl(1-q_i^{-em_i}\bigr)^{-1}-1\Bigr],\qquad(k_i)=(2,2,1,1).
\]
The second bracket bounds all nonempty changes of high digits by geometric series, since
ternary subset sums have distinct exponents. The bound decreases when a prime increases, and
for each base tuple the program checks $1-1/\lambda-T(b)>1/200$ in rational arithmetic.
Therefore
\[
 W_q(C)-W_q(F)\ \ge\ M\bigl(1-1/\lambda-T(q)\bigr)\ \ge\ 0. \qedhere
\]
\end{proof}

Call a rational ratio \emph{$P$-supported} if every prime factor of its numerator and
denominator lies in $P$, and let $R_P$ be the set of positive integers $m$ with
$v_p(m)\in\mathcal A$ for all $p\in P$, with no condition at the other primes. Since
$\mathcal A$ is progression-free, $R_P$ avoids all progressions with $P$-supported ratio.

\begin{corollary}[Restricted ratios; \stU]\label{gpi:cor:restricted}
Let $P$ be a finite set of primes on which Rankin's rule is optimal, for instance any set of
at most four primes. Then $R_P$ attains the largest natural density, and the largest upper
logarithmic density, among the sets of positive integers avoiding all progressions with
$P$-supported rational ratio. Its density is
\[
 d_P=\prod_{p\in P}\Bigl[(1-p^{-1})\prod_{j\ge0}(1+p^{-3^j})\Bigr].
\]
In particular $\alpha\le d_{\set{2,3,5,7}}<0.740912$.
\end{corollary}

\begin{proof}
The set $R_P$ has natural density $d_P=c_PW(C_P)$ by Lemma~\ref{gpi:lem:product}, with
$A_p=\Z_{\ge0}$ outside $P$. Let $E$ avoid all progressions with $P$-supported ratio and write
$m=hs$ with $s\in S_P$ and $h$ coprime to the primes of $P$. For each $h$ the fiber
$\set{s:hs\in E}$ is progression-free, so its reciprocal sum is at most $W(C_P)$, and
\[
 \sum_{m\in E,\ m\le x}\frac1m\ \le\ W(C_P)\sum_{h\le x}\frac1h=(c_PW(C_P)+o(1))\log x
\]
by periodicity of the coprimality condition and partial summation. Hence the upper
logarithmic density of $E$ is at most $d_P$, and so is its natural density if it exists.
Sets free of all rational-ratio progressions obey the restriction for every $P$, which gives
the bound for $\alpha$.
\end{proof}

The values for the three initial segments of the primes are
\begin{center}\small
\begin{tabular}{lr}
\toprule
$P$ & $d_P$ (rounded)\\ \midrule
$\set{2,3}$ & $0.779336473$\\
$\set{2,3,5}$ & $0.754148705$\\
$\set{2,3,5,7}$ & $0.740911746$\\ \bottomrule
\end{tabular}
\end{center}
For these three sets the fixed-prime certificates of
\gpifile{mcnew_density_and_grh_transfers/checks/verify_rankin.py}
(low boxes of sizes $81$,
$243$, $729$ with $\lambda=51/50$, $11/10$, $21/20$) give
$W(C_P)-W(F)\ge\theta\,W(C_P\setminus F)$ with $\theta>0.014$, $0.062$, $0.0057$
respectively, and hence uniqueness of the optimizer on the smooth semigroup. Each of the
three sets lies outside the range $1/c_P\le2$ of Proposition~\ref{gpi:prop:light}.

\subsection{How many primes must a density rule couple?}\label{gpi:sec:blocks}
{\raggedright\sources{grh_prime_patterns_and_branch_precision; mcnew_counterexample_and_grh_transfers; prime_observations_and_progression_free_sets}}

Fix $k\ge3$ and $d\ge1$. Partition the primes into blocks of size at most $d$, choose for each block
$P$ an arbitrary set $F_P\subseteq S_P$ of allowed smooth parts, and admit a positive integer when its
$P$-smooth part lies in $F_P$ for every block. Let $\alpha_{k,d}$ be the supremum of the upper
densities of the sets of this form that are free of $k$-progressions, over all partitions and all
rules, and let $\alpha_k$ be the supremum of the natural densities of all sets free of
$k$-progressions; thus $\alpha_3=\alpha$. The partition may vary with the competitor, so the model
allows dependence among infinitely many separate groups of primes. A nonempty progression-free
product has progression-free block rules (replace the smooth part of one admitted integer in one
block by the terms of a putative local progression), and conversely a product of progression-free
rules is progression-free, since in a progression of the product the smooth parts in each block form
a progression with the block part of the ratio. Write $p_j$ for the $j$th prime.

\begin{theorem}[Bounded prime blocks at every length; \stU]\label{gpi:thm:blocksk}
For all $k\ge3$ and $d\ge1$,
\[
 \alpha_{k,1}=d(\mathcal G_k),\qquad
 0<\alpha_k-\alpha_{k,d}\le\frac{(p_{d+1}-1)^{-(k-2)}}{k-2}.
\]
Every $\alpha_{k,d}$ is attained by a set with a natural density that uses the greedy rule $G_k$ off
finitely many primes, and the numbers $\alpha_{k,d}$ and $\alpha_k$ are computable, uniformly in
$k$ and $d$.
\end{theorem}

Thus every fixed bound on the number of jointly tested primes has a strict density cost, even after
optimizing the partition.

\begin{proof}
\emph{Local probabilities and width one.} For a finite set $P$ of primes and $F\subseteq S_P$, the
integers with $P$-smooth part in $F$ have density $c_PW(F)$ by \cref{gpi:lem:product}, and disjoint
finite blocks are independent. By \cref{gpi:prop:light}, every block $P$ with $1/c_P\le2$, in
particular every singleton, satisfies $W(F)\le W(G_k^P)$ for progression-free $F$. Finite collections
of block restrictions bound the upper density of a width-one competitor by finite products of greedy
factors, which converge to $d(\mathcal G_k)$; since $\mathcal G_k$ has width one,
$\alpha_{k,1}=d(\mathcal G_k)$.

\emph{Reduction to finitely many exceptional primes.} Call a block \emph{heavy} if $1/c_P>2$. If all
primes of $P$ exceed $4d$, then $\log(1/c_P)\le\sum_{p\in P}1/(p-1)\le\tfrac14<\log2$; so at most
$\pi(4d)$ blocks are heavy, and their union $H$ has at most $h=d\,\pi(4d)$ primes. Replace every
other block by the greedy rule. The new set has a natural density, equal to the finite heavy
probability times the greedy product off $H$: ignoring the primes $p>Y$ changes at most
$x\sum_{p>Y}p^{-(k-1)}$ integers up to $x$, since $0,\dots,k-2\in G_k$. By the weight comparison it
dominates the upper density of the original set, which is bounded by the probability of passing any
finite collection of the original blocks.

\emph{Attainment.} For fixed finite $P$ the maximum of $W$ over progression-free $F\subseteq S_P$ is
attained, by a pointwise convergent subsequence of indicators on the countable exponent lattice.
Start with a maximizing sequence in the reduced form and pass to a subsequence along which the
number of exceptional primes and the partition of their ordered positions are fixed, and each prime
label is either constant or tends to infinity. Imposing valuation zero at the escaping labels costs
at most the sum of their reciprocals, which tends to zero; the resulting slices on the fixed labels
are progression-free. Replace each by an optimal rule and use $G_k$ at every other prime, which is
possible because $0\in G_k$; this single construction dominates the limiting densities.

\emph{Strict loss.} Take an optimizer of this form. Off its finitely many exceptional primes,
\cref{gpi:thm:klength} improves the greedy rule on a disjoint finite set of primes. The unchanged
block factors have positive product, since the optimizer dominates $\mathcal G_k$, so the
improvement gains a positive amount of natural density, and $\alpha_k>\alpha_{k,d}$.

\emph{Upper bound for the loss.} Put $Y=p_{d+1}-1$, $P=\set{p\le Y}$ (a set of $d$ primes) and
$M=\max W(F)$ over progression-free $F\subseteq S_P$. In any progression-free set $A$ the integers
with a fixed $P$-rough part $u$ have a progression-free set of smooth parts, so
\[
 \sum_{n\le x,\ n\in A}\frac1n\le M\sum_{\substack{u\le x\\(u,\prod_{p\in P}p)=1}}\frac1u=Mc_P\log x+O_P(M),
\]
and partial summation gives $d(A)\le Mc_P$ whenever $d(A)$ exists. An optimizing block on $P$ with
greedy singletons outside it gives $\alpha_{k,d}\ge Mc_P\prod_{p>Y}(1-p^{-(k-1)})$, and $Mc_P\le1$. The
bound follows from $1-\prod_j(1-x_j)\le\sum_jx_j$ and $\sum_{j>Y}j^{-(k-1)}\le Y^{-(k-2)}/(k-2)$.

\emph{Computation.} Freezing the exceptional primes above $Z$ to valuation zero costs at most $h/Z$,
and releasing them to the greedy rule can only increase the density. Enumerate all partitions of
the primes up to $Z$ into blocks of size at most $d$; in each block truncate the exponents below
$L$, enumerate all subsets of the finite box and keep the progression-free ones. Extending a
truncated rule by excluding higher exponents preserves admissibility and loses at most
$\sum_{p\le Z}p^{-L}$, and the greedy factors on the remaining primes can be evaluated to any certified
accuracy by the same truncation and the bound $\sum_{p>U}p^{-(k-1)}$. These explicit errors give a
uniform algorithm for $\alpha_{k,d}$, and the bound for the loss gives one for $\alpha_k$.
\end{proof}

For $k=3$ write $\alpha_m=\alpha_{3,m}$. By \cref{gpi:thm:blocksk} it is also the supremum of the
natural densities of the progression-free block sets of width $m$ that have a natural density, it is
attained, $d(R)\le\alpha_1\le\alpha_2\le\dots<\alpha$, and $\alpha-\alpha_m\le1/(p_{m+1}-1)$, so
that $\alpha_m\to\alpha$. The constant and the small widths can be sharpened.

\begin{theorem}[Bounded prime blocks for three-term progressions; \stU]\label{gpi:thm:blocks}
For every $m\ge1$,
\[
 0<\alpha-\alpha_m\le\frac1{p_{m+1}},\qquad\text{so that}\quad\lim_{m\to\infty}\alpha_m=\alpha ,
\]
and
\[
 \alpha_1=\alpha_2=\alpha_3=\alpha_4=d(R).
\]
The least $m$ with $\alpha_m>d(R)$ satisfies $5\le m\le225$. The positivity of $\alpha-\alpha_m$
and the limit are the case $k=3$ of \cref{gpi:thm:blocksk}; the new content is the constant
$1/p_{m+1}$, the equalities for $m\le4$ and the range for the least $m$.
\end{theorem}

A progression-free product rule of density greater than $d(R)$ must therefore couple at least five
primes in some block.

\begin{proof}
\emph{The constant.} Positivity and the limit are \cref{gpi:thm:blocksk} with $k=3$. Take
$Y=p_{m+1}-1$, $P$ and $M$ as in the upper bound for the loss in its proof, which gives
$\alpha\le c_PM\le1$ and $\alpha_m\ge c_PM\prod_{p>Y}(1-p^{-2})$. Telescoping over all integers
greater than $Y$ gives $\prod_{p>Y}(1-p^{-2})\ge Y/(Y+1)$, and hence
$\alpha-\alpha_m\le1/(Y+1)=1/p_{m+1}$.

\emph{Small widths.} Rankin's rule is itself a rule of width one, so $\alpha_m\ge d(R)$. By
\cref{gpi:thm:fourprimes}, Rankin's rule is optimal on every block of a rule of width $m\le4$, so the
replacement argument of \cref{gpi:thm:blocksk} applies to all blocks, heavy or light, and gives
$\alpha_m\le d(R)$ for $m\le4$. The set of \cref{gpi:thm:explicit} is a rule of width $225$ of density
greater than $d(R)$.
\end{proof}

In particular, for a fixed finite set $T$ of primes, no progression-free rule on $S_T$ combined with
Rankin's rule at all other primes attains $\alpha$: such a set has density at most
$\alpha_{\#T}<\alpha$. Exact optimality on finite sets of primes, as in \cref{gpi:thm:fourprimes},
therefore coexists with the failure of the unrestricted conjecture.

\subsection{Integer ratios against rational ratios: Question 7.4}\label{gpi:sec:q74}
{\raggedright\sources{grh_transfer_extensions_2026-11-26; mcnew_density_and_grh_transfers}}

McNew constructed a set avoiding integer ratios that is denser than Rankin's set
\cite[Theorem~3.3]{McNew2015Sets}. Comparing one construction with Rankin's set leaves open
whether another set avoiding rational ratios is equally dense; Question~7.4 asks for the
comparison of the two suprema. The following comparison is uniform in the competing set.

\begin{theorem}[Question 7.4; \stU]\label{gpi:thm:beta-alpha}
\[
 \beta-\alpha\ \ge\ \frac{135102729954660}{10^{18}}>0.0001351 .
\]
\end{theorem}

The proof replaces the $5$-smooth part of an arbitrary competitor by one fixed finite set
and pays for the new progressions that cross the boundary of the replacement. Let
$S=S_5=\set{2^i3^j5^k:i,j,k\ge0}$ and
\[
 N=10^6,\qquad\tau=\frac{824341}{10^7},\qquad g=\frac{237}{100000},\qquad V=S\cap[1,N].
\]
The set $V$ has $507$ elements. Let $B$ be the set of $102$ integers in
\cref{gpi:tab:replacement}. It is obtained as follows: start with the seven integers
$60,90,150,180,300,450,900$ added to the slice at $2,3,5$ in the proof of McNew's
Theorem~3.3; then visit the integers $2^i3^j5^k\le N$ with all ternary digits of $i,j,k$ in
$\set{0,1}$ in increasing order, retaining an integer if adjoining it introduces no
integer-ratio progression.

\begin{table}[htbp]\centering\small
\caption{The replacement set $B$.}\label{gpi:tab:replacement}
\begin{tabular}{rrrrrrrrrr}\toprule
1&2&3&5&6&8&10&15&16&24\\
27&40&48&54&60&80&81&90&120&125\\
135&150&162&180&216&250&270&300&375&405\\
432&450&512&625&648&750&900&1000&1024&1080\\
1250&1296&1536&1875&2000&2560&3000&3072&3375&4096\\
5000&5120&6750&7680&8192&10000&10125&12288&13824&15360\\
16875&19683&20480&24576&27000&27648&39366&40960&41472&50625\\
59049&61440&64000&69120&82944&98415&110592&118098&122880&128000\\
138240&157464&192000&196830&207360&221184&295245&314928&320000&331776\\
384000&414720&472392&512000&531441&552960&590490&640000&663552&787320\\
944784&960000&&&&&&&&\\\bottomrule
\end{tabular}
\end{table}

\begin{lemma}[Finite exchange with a boundary charge]\label{gpi:lem:exchange}
The set $B$ contains no integer-ratio progression. Every $A\subseteq S$ free of
rational-ratio progressions satisfies
\begin{equation}\label{gpi:eq:exchange}
 W(B)-W(A)-\tau\,W(B\setminus A)\ \ge\ g .
\end{equation}
Moreover $\sum_{r\in\N\setminus S}r^{-2}<\tau$.
\end{lemma}

\begin{proof}
Checking, for every $a\in B$ and every integer $2\le r\le\sqrt{N/a}$, that $ar$ and $ar^2$
are not both in $B$ exhausts the integer-ratio triples in $B$. Geometric series give
$W(S)=15/4$, and exact summation gives
\[
 W(B)=\frac{1928068463873}{680244480000},\qquad
 T_0:=W(S\setminus V)=\frac{12000753838363}{108839116800000000}.
\]
Give $m\in V$ the weight $\omega_m=(1-\tau\mathbf 1_B(m))/m$. The finite statement to be
certified is
\begin{equation}\label{gpi:eq:finite}
 \sum_{m\in I}\omega_m\ \le\ (1-\tau)W(B)-T_0-g
 \qquad\text{for every }I\subseteq V\text{ free of rational-ratio progressions.}
\end{equation}
Since $W(A\setminus V)\le T_0$ and $B\subseteq V$, applying \eqref{gpi:eq:finite} to
$I=A\cap V$ and rearranging gives \eqref{gpi:eq:exchange}.

All weights are scaled to integers by the factor $1088391168000000000000000$. In these units
the bound in \eqref{gpi:eq:finite} is
\[
 2827908305897851139120000 ;
\]
equivalently, every set of vertices meeting each rational-ratio triple in $V$ has weight at
least $999136824871360000000000$.
There are $16{,}446$ such triples, generated from the pairs $a<c$ in $V$ with $ac$ a square
and $\sqrt{ac}\in V$. Two independent verifications are provided by
\gpifile[.]{grh_transfer_extensions_2026-11-26/code/certify_density.py}
\begin{enumerate}
\item A proof tree for the covering bound, with two kinds of steps. At a leaf, nonnegative
 integer masses are assigned to the remaining triples so that the total mass incident to each
 vertex is at most its weight; every cover then costs at least the total mass, and the leaf
 is certified when this total reaches the deletion weight still required. Otherwise one
 uncovered triple is chosen; every cover contains one of its three vertices, and one branches
 on the three possibilities, paying the weight of the chosen vertex and removing every triple
 it meets. The tree, with every branch and the dual masses of every leaf, is recorded in
 \gpifile[;]{grh_transfer_extensions_2026-11-26/data/density_exchange_certificate.json}
 a separate replayer checks the capacities, the uncovered-edge conditions and all three
 children at each branch. Generation and replay each visit $496$ nodes.
\item A direct search. Vertices are visited in increasing order; whenever a middle term $y$
 and an earlier selected term $x$ are present, the future term $y^2/x$ is excluded if it is an
 integer in $V$. Every admissible set follows one of the include/exclude branches, and the
 sum of the selected and still available weights bounds every completion. The search closes
 every branch at or below the bound in $416{,}653$ visits.
\end{enumerate}
Finally $\sum_{r\notin S}r^{-2}=\zeta(2)-\prod_{p\in\set{2,3,5}}(1-p^{-2})^{-1}=\zeta(2)-25/16$.
Bounding the tail of $\sum r^{-2}$ beyond $10000$ by $1/10000$ and rounding each earlier term
upward on the grid $10^{-18}\Z$ gives
\[
 \zeta(2)-\frac{25}{16}<\frac{82434071848064747}{10^{18}}<\frac{824341}{10^7}. \qedhere
\]
\end{proof}

\begin{lemma}[The exchange survives further primes]\label{gpi:lem:patch}
Let $Y\ge5$. For every $E\subseteq S_Y$ free of rational-ratio progressions there is an
$F\subseteq S_Y$ free of integer-ratio progressions with $W(F)\ge W(E)+g$.
\end{lemma}

\begin{proof}
Put $A=E\cap S$ and replace $A$ by $B$, keeping the rest of $E$. For each newly inserted
$b\in B\setminus A$ delete every $br^2$ in the current set for which the integer $r$ has a
prime factor greater than $5$; these deletions lie outside $S$. Let $F$ be the result.

Consider an integer-ratio triple in $F$ that was not in $E$; it contains a newly inserted
element of $B$. If its ratio is $5$-smooth, the valuations at the primes greater than $5$ are
constant along it, so it lies entirely in $S\cap F=B$, where there is no such triple. If its
ratio has a prime factor greater than $5$, the middle and last terms are not in $S$, so the
inserted element is the first term, and the last term was deleted. A triple consisting of
old elements only was already forbidden in $E$.

The deleted mass is at most $W(B\setminus A)\sum_{r\notin S}r^{-2}<\tau\,W(B\setminus A)$,
even when all positive integer ratios are allowed for, uniformly in $Y$ and in $E$. Now apply
Lemma~\ref{gpi:lem:exchange}.
\end{proof}

The uniformity in $Y$ is essential: an improvement on a fixed finite set of primes alone
would have no consequence for natural densities. The ordering imposed by integer ratios
(nonnegative steps in all coordinates) is what allows a new cross-boundary triple to be
charged to its final term; rational ratios allow mixed signs.

\begin{proof}[Proof of \cref{gpi:thm:beta-alpha}]
Fix $Y\ge5$ and use Lemma~\ref{gpi:lem:BU} for $K=\Q$ and both conventions, writing $B_{\Q}$ and
$B_{\mathrm I}$ for the two values of the weighted optimum $B$. By Lemma~\ref{gpi:lem:patch},
$B_{\mathrm I}\ge B_{\Q}+cg$. Hence
\begin{equation}\label{gpi:eq:bridge}
 \beta-\alpha\ \ge\ vB_{\mathrm I}-B_{\Q}\ \ge\ cg-(1-v)B_{\mathrm I}\ \ge\
 g\prod_{p\le Y}(1-p^{-1})-\sum_{p>Y}p^{-2}.
\end{equation}
Take $Y=6619$ and sieve the primes up to $2{,}000{,}000$. Starting from the integer
$10^{18}$, multiply successively by $(p-1)/p$ and round downward for each $p\le Y$; for
$Y<p\le2{,}000{,}000$ round $10^{18}/p^2$ upward, and add $10^{18}/2{,}000{,}000$ to dominate
the remaining tail. This proves
\[
 \prod_{p\le Y}(1-p^{-1})\ge\frac{63728234915742846}{10^{18}},\qquad
 \sum_{p>Y}p^{-2}\le\frac{15933186795650}{10^{18}},
\]
and substitution into \eqref{gpi:eq:bridge}, with one final downward rounding, gives the
theorem.
\end{proof}

Even the elementary bounds $\prod_{p\le Y}(1-p^{-1})\ge1/(2\sqrt Y)$ and
$\sum_{p>Y}p^{-2}\le1/Y$ make the right side of \eqref{gpi:eq:bridge} positive for large $Y$;
the first follows by including every odd integer $3\le j\le Y$ in the product and using
$(1-1/j)^2\ge(j-2)/j$.

\subsection{Integral ratios in number fields: a larger exchange}\label{gpi:sec:cube}
{\raggedright\sources{grh_ideal_densities_and_transfer_conditions_2026-11-26; mcnew_counterexample_and_grh_transfers}}

For integral ratios the improvement on the greedy density guaranteed by
\cref{gpi:thm:onedigit} can be made much larger by an exchange that removes a whole orthant.
Its integer precedent is McNew's exchange at $2,3,5$, and its polynomial analogue is
\cref{gpp:thm:natural}. For a nonempty finite set $S$ of prime ideals of $K$, with $k=\#S$,
replace the permitted local vectors $\mathcal A^k$ by the set of \cref{gpp:sec:natural},
\[
 D_S=\set{u\in\mathcal A^k:\min_iu_i=0}\ \cup\ \bigl(\set{1,2}^k\setminus\set{(1,\dots,1)}\bigr),
\]
and let $\mathcal E_S$ be the set of ideals whose exponent vector on $S$ lies in $D_S$ and
whose exponents outside $S$ lie in $\mathcal A$. With $T(w)$ as in \cref{gpp:thm:natural}, put
\[
 C_S=\prod_{\mathfrak p\notin S}\Bigl[(1-w_{\mathfrak p})\sum_{a\in\mathcal A}w_{\mathfrak p}^a\Bigr].
\]

\begin{proposition}[Orthant exchange; \stU]\label{gpi:prop:cube}
The set $\mathcal E_S$ is admissible for the integral convention, has a natural density, and
\begin{equation}\label{gpi:eq:cube-delta}
 d(\mathcal E_S)-d(\mathcal G_K)=C_S\prod_{\mathfrak p\in S}(1-w_{\mathfrak p})w_{\mathfrak p}
 \Bigl[\prod_{\mathfrak p\in S}(1+w_{\mathfrak p})-1-\prod_{\mathfrak p\in S}T(w_{\mathfrak p})\Bigr].
\end{equation}
Moreover $\prod_{\mathfrak p\in S}T(w_{\mathfrak p})<2^n$ and $C_S>2^{-n}$, so the gain is
positive as soon as
\begin{equation}\label{gpi:eq:cube-test}
 \prod_{\mathfrak p\in S}(1+w_{\mathfrak p})>1+2^n ,
\end{equation}
which holds for all sufficiently large initial segments $S$ of the primes of $K$. Replacing
$C_S$ by $2^{-n}$ and $\prod T(w_{\mathfrak p})$ by $2^n$ in \eqref{gpi:eq:cube-delta} then
gives a positive rational lower bound for the gain.
\end{proposition}

\begin{proof}
A forbidden triple in $\mathcal E_S$ has zero valuation step outside $S$, since $\mathcal A$
is progression-free. Inside $S$ it would give $u,u+h,u+2h\in D_S$ with $h\ge0$, $h\ne0$, which
\cref{gpp:lem:ES} excludes.

The densities exist by Lemma~\ref{gpi:lem:product}. The deleted vectors are those of
$\mathcal A^k$ with all coordinates positive, of mass
$\prod_{\mathfrak p\in S}w_{\mathfrak p}T(w_{\mathfrak p})$, and the added cube has mass
$\prod_{\mathfrak p\in S}w_{\mathfrak p}\,[\prod_{\mathfrak p\in S}(1+w_{\mathfrak p})-1]$; this gives
\eqref{gpi:eq:cube-delta}. The exponents occurring in $T$ lie in $\set0\cup\set{2,3,\dots}$,
and every integer at least two is $2a+3b$ with $a,b\ge0$, whence
\[
 T(w)\le\frac1{(1-w^2)(1-w^3)},\qquad
 \prod_{\mathfrak p\in S}T(w_{\mathfrak p})\le\zeta_K(2)\zeta_K(3)\le\bigl(\zeta(2)\zeta(3)\bigr)^n<2^n .
\]
The middle inequality holds prime by prime, since at most $n$ prime ideals lie over a
rational prime and each has norm at least that prime; the last follows from
$\zeta(2)<\sum_{j\le10}j^{-2}+\tfrac1{10}<\tfrac{33}{20}$ and
$\zeta(3)<\sum_{j\le4}j^{-3}+\tfrac1{32}<\tfrac{121}{100}$, whose product is $3993/2000$.
Also $C_S\ge\zeta_K(2)^{-1}>2^{-n}$ because $0,1\in\mathcal A$. Finally
$\prod_{\mathfrak p\in S}(1+w_{\mathfrak p})$ is unbounded along increasing $S$, because
$\sum_{\mathfrak p}w_{\mathfrak p}$ diverges.
\end{proof}

Enumerating and factoring rational primes gives a terminating search for $S$, and factoring
an ideal and testing its valuation digits decides membership in $\mathcal E_S$. The exchange
does not serve for fractional ratios, by the mixed-sign progression of \cref{gpp:rem:mixed};
for that convention the improvement is \cref{gpi:thm:onedigit}.

\begin{proposition}[An explicit gain in $\Q(\sqrt{-23})$; \stU]\label{gpi:prop:Q23gain}
Let $K=\Q(\sqrt{-23})$ and let $S$ consist of the four prime ideals of norms $2,2,3,3$. Then,
for the integral convention,
\[
 d(\mathcal E_S)-d(\mathcal G_K)>\frac{16658080679}{13931406950400}>0.0011957 .
\]
\end{proposition}

\begin{proof}
Write $\omega=(1+\sqrt{-23})/2$, with minimal polynomial $X^2-X+6$. For $p=2,3$ its two
distinct roots modulo $p$ are $0$ and $1$, so the four prime ideals $(p,\omega)$ and
$(p,\omega-1)$ have norms $2,2,3,3$. The bound used in the proof of \cref{gpp:thm:natural},
which follows from $\mathcal A\cap[1,8]=\set{1,3,4}$, is
\[
 T(w)\le\overline T(w):=1+w^2+w^3+\frac{w^8}{1-w},\qquad
 \overline T(1/2)=\frac{177}{128},\quad\overline T(1/3)=\frac{5023}{4374},
\]
and exact arithmetic gives
\[
 \prod_{\mathfrak p\in S}(1+w_{\mathfrak p})-1-\prod_{\mathfrak p\in S}\overline T(w_{\mathfrak p})
 =\frac{16658080679}{34828517376},\qquad
 \prod_{\mathfrak p\in S}(1-w_{\mathfrak p})w_{\mathfrak p}=\frac1{324}.
\]
Since $\prod_{\mathfrak p\in S}(1-w_{\mathfrak p}^2)=4/9$ and
$\zeta_K(2)\le\zeta(2)^2<(5/3)^2$,
\[
 C_S\ \ge\ \prod_{\mathfrak p\notin S}(1-w_{\mathfrak p}^2)=\frac{\zeta_K(2)^{-1}}{4/9}>\frac{81}{100}.
\]
Substituting into \eqref{gpi:eq:cube-delta} gives the bound. The prime reductions and every
number in this certificate are checked by
\gpifile[.]{grh_ideal_densities_and_transfer_conditions_2026-11-26/code/verify_ideal_natural_gain.py}
\end{proof}

\subsection{Open questions}\label{gpi:sec:open}
{\raggedright\sources{grh_transfers_density_gaps_and_zero_tests; mcnew-density-and-dirichlet-grh; mcnew_counterexample_and_grh_transfers; mcnew_density_and_grh_transfers; grh_prime_patterns_and_branch_precision; grh_ideal_densities_and_transfer_conditions_2026-11-26; prime_observations_and_progression_free_sets}}

\begin{enumerate}
\item \emph{The values of the four suprema.} The results above give
 \[
 d(R)+10^{-596}<\alpha<0.740912,\quad0.723891<\beta<0.749192,\quad
 0.7365954<\overline\alpha<0.7716011,
 \]
 together with $\beta-\alpha>0.0001351$ and $\overline\beta>0.818410$. The numbers $\alpha$,
 $\beta$ (and $b_K$, $u_K$ in general) are computable, but none of them is determined. How
 large is $\alpha-d(R)$? The exchanges used here optimize particular families of patterns;
 larger digit alphabets or several interacting markers could give larger gains.
\item \emph{The least number of coupled primes.} Determine the least $m$ with
 $\alpha_m>d(R)$; \cref{gpi:thm:blocks} gives $5\le m\le225$. More generally, for each $k\ge3$, what
 is the least $d$ with $\alpha_{k,d}>d(\mathcal G_k)$ in \cref{gpi:thm:blocksk}? In particular, is
 Rankin's rule optimal on every set of five primes? Can the positive margin in the low-box inequality
 \eqref{gpi:eq:lowbox} be certified on larger sets of small primes without a low-digit box
 that grows exponentially?
\item \emph{More deletions.} \Cref{gpi:thm:edge} classifies the profitable one-digit
 exchanges with at most two deleted cube vertices. Is there a comparably explicit
 classification for three or more deletions, or an exchange with an appreciable density gain
 using few primes of ordinary size?
\item \emph{Visible gains for longer progressions.} The gain of \cref{gpi:thm:klength} can be
 extremely small. Can a tractable finite witness visibly improve on the greedy density
 $0.89537\ldots$ at length four?
\item \emph{The loss at each prime.} Is $\overline\alpha_{p^-}-\overline\alpha_p>0$ for every
 sufficiently large prime $p$, where $p^-$ is the preceding prime? The proof at $p=11$ in
 \cref{gpi:thm:upper-alpha} reduces strictness to finding one interval on which the
 independent layers have a strictly larger optimum; a family of such intervals would
 quantify the convergence from smooth to unrestricted ratios.
\item \emph{The size of $\beta-\alpha$.} Can the replacement set $B$ and the budget for the
 removed endpoints in Lemmas~\ref{gpi:lem:exchange} and~\ref{gpi:lem:patch} be optimized jointly to improve
 the lower bound of \cref{gpi:thm:beta-alpha}?
\item \emph{Number fields.} Determine the optimal natural density $b_K$ and the extremal upper
 density $u_K$, or sharp field-specific improvements on the greedy density, for instance for
 quadratic fields as asked in \cite{Best2017Geometric}. For fractional ratios the
 improvements obtained here are those of \cref{gpi:thm:onedigit}.
\end{enumerate}
